% Planetary Research - A diamond open access journal
% Manuscript built on the journal's LaTeX template (journal-templates, main branch, template dated July 18, 2026).
% Build: pdflatex Hollow_Formation_Model_Paper_final.tex && biber Hollow_Formation_Model_Paper_final && pdflatex (twice). Bibliography: references_final_rumil.bib (written by Rumil check_latex).
% Version 3 (2026-10-30): addresses the author's comments on v2 (see response_to_comments_v3.md), merges the two
% Caloris nodes of the crater containment graph (analysis_v8), and revises the charging and collapse steps.
% Version 4 (2026-11-01): step 2 becomes "Charging and collapse", with collapse while the crater is hot kept in the
% main text and compared with karst and thermokarst; trims from the evidence audit (cut text in cut_material_v4.tex);
% crater-graph numbers from analysis_v9 (peak-ring uplift beneath peak-ring craters and basins); a mantle source
% for the deepest hosts; supplement references follow supplementary-materials_v9 (S9 slope cut, old S12 merged into
% S9, old S15 merged into S5).
% Version 5 (2026-11-02): applies the author's 20 comments on the v4 PDF (LRM link softened, mantle source, AIC
% explained, Figure 5 timescales, plains hosts, fields by scarp retreat, wording) and the hollow-model workshop:
% craters with explosive vents (new Section 3.3; analysis_workshop/test8_vents.py), depth quantiles and sunlight
% (analysis_workshop/test1_depth_quantiles.py, test1b_additive_solar.py, test2_depth_size.py), two routes to a uniform
% base temperature (steady conduction, sulfur heat pipe), depth formed after exposure as a disfavored alternative,
% floor renewal and halo relief added to the untested steps and predictions. Built by _scratch_v10_make_v5.py;
% supplement: supplementary-materials_v10.tex.
% Version 6 (2026-11-03): applies the author's 50 comments on the v5 PDF: Section 3.4 rewritten (hot longitudes at
% +-15 deg, crater population, continuous rise with Delta T_w, Stage 1 freshness; analysis_workshop/insolation_v6),
% depth vs low-reflectance material and the X-ray caveats (analysis_workshop/composition_v6), the volatile comparison
% (analysis_workshop/volatile_fit), a joint model of hosting (new Section 3.7; analysis_workshop/summary_model), the
% impact pane in Figure 5, and the Caloris test on the v9 graph (analysis_workshop/caloris_v9). Built by
% _scratch_v11_make_v6.py; supplement: supplementary-materials_v11.tex.
% Color (2026-11-03): MDIS eight-filter map results in Sections 2.3, 3.6, 3.7, and 4.2 (analysis_workshop/color_mdr,
% summary_model/step9_color.py, color_mdr/depth_color_lrm.py); Figure S7 is the new color figure, so the T_base figure is S8.
% Version 6.4 (2026-11-07): Barraud et al. (2023) is no longer cited. The sulfide and Cl contents of the plains come from
% McCoy et al. (2018), so Section 4.1 now states that sulfide content does not explain the scarcity of hollows on the
% plains, and the chloride sentence in Section 4.3 adds the Cl of the Northern Terrane and the collapse temperatures of
% Section 4.2. Numbers checked against McCoy et al. (2018) Table 7.1 (normative sulfides, wt%: Caloris interior plains
% 3.7, Low-Fast 3.7, Northern 4.3, southern hemisphere 4.7, high-Mg 6.8; Cl 0.15, 0.24, 0.45, 0.14, 0.13).
% Version 6.5 (2026-11-07): the screened v3 catalog with the refit detection model (code/rerun/out_v3; numbers from
% rerun/changes_all_v3r_g7.csv, g4 rows for completeness-dependent values) in every section and figure; terrain classes traced
% by hand in 64 hosts (rerun/g13_traced, revision r2) and curvature classes elsewhere; the model rewrite of
% model_rewrite_draft_v2 with the author's annotations on its PDF (silica reduction, enclosed impact melt, peak height and
% floor position, the knobs around Caloris, the floating sulfide layer, prediction 9). Supplement: supplementary-materials_v14.tex.
% Version 6.6 (2026-11-07): terrain results from the 62 host craters traced by hand alone (rerun/g16_traced_only), without
% automated (curvature) classes; claims that fail with the traced hosts removed (hollow fields high on peaks, hollows densest at floor
% centers); the author's notes on Section 3.8 (v6_5_markup.pdf); cuts and condensations agreed on 7 Nov (368 K layer, simple-to-complex
% step, fill-only result, knob comparison with Class 1-2 peaks, knob color and depths, Sections 3.6, 3.7, 4.1, 4.2, Conclusions).
% Supplement: supplementary-materials_v15.tex.
% Version 6.7 (2026-11-07): terrain results from the 111 host craters traced in rounds 1 and 2 (all hosts 30-250 km with >= 5 hollows;
% rerun/g18_traced_r3). Supplement: supplementary-materials_v16.tex.
% Version 6.8 (2026-11-07): knobs compared with the peaks of Class 1-2 craters traced by hand (rerun/g19_knobs_traced).
% Supplement: supplementary-materials_v17.tex.
% Version 6.9 (2026-11-09): Section 3.8 locates the regional clustering of hosting; new Figure 6 (regional excess over the joint model;
% rerun/g20_regional_clusters). The life-cycle figure becomes Figure 7. Supplement: supplementary-materials_v19.tex.
% Version 6.10 (2026-11-09): Michael's strikeout on v6.8 p. 12 (34 K step sentence) applied; sweep for redundancy (sweep_v6_10_v20_changes.md).
% Version 6.11 (2026-11-09): Michael's 92 comments on the v6.10 PDF; two-panel Figure 1; hollow depth from the depths of
% Blewett et al. (2016) and from new automatic and hand shadow depths, with depth against hollow size (rerun/g23, g24; new Figure 5);
% composition group in the joint model and the regional excess refit (rerun/g25; new Figure 6); the province screen; life-cycle
% figure after the six steps with the thermal panel as an inset (Figure 7); Figure S8 promoted to Figure 8 (Section 4.2).
% Version 6.12 (2026-11-09): measurement uncertainty of the shadow depths, briefly in Sections 3.7 and 4.2 and Figure 8 (details in
% supplement S9; rerun/g26); Section 2.3 sentence split (Michael).
% Version 6.13 final (2026-11-09), in submission/: Michael's comments on v6.12 (Section S9 pointer moved in 2.3; "The automatic
% depths"; Figure 5 legend moved, fig_depth_v3; latitude paragraph of 4.2 and the limitations paragraph of 4.4 removed; Figure 8 cited
% in 4.2; prediction 1 wording). Supplement: supplementary-materials_v24_final.tex (same content as v24).
% Version 6.14 final (2026-11-09), in submission/: numbers corrected to the final (v3r) completeness grid: 3.3 crater counts with
% completeness >= 0.5 (1,252; 414); 3.2 Caloris comparisons (2 of 840; 1 of 58) and Figure 2c from the v3r rerun of the Caloris tests
% (rerun/g3_catalog/caloris/out_v3r; fig1_exposure_v9_v3cat_traced_r3_v3r). Supplement: supplementary-materials_v25_final.tex.
% Version 6.15 (2026-11-10), submission/Hollow_Formation_Model_Paper_final.tex: Figure 2a on the MDIS mosaic with the coverage shown by
% pixel scale (rerun/g28_fig2a; fig1_exposure_v10); Figure 7 lettered a-g, with the step numbers in circles and a deeper collapse
% (figscripts/lifecycle_fig_v12.py), its caption, and two references to its panels (Michael). Supplement: supplementary-materials_final.tex (v26).
% Bibliography: references_final_rumil.bib (Rumil check_latex).
\documentclass[article,review]{planetary}
\usepackage{placeins}

% ------ Title block --------------------------------------------------
\title{Toward a Geologically Plausible Model of Hollow Formation on Mercury}
\shorttitle{Toward a geologically plausible model of hollow formation on Mercury}

% ------ Author block -------------------------------------------------
\author[1]{Michael S. Phillips \orcid{0000-0001-8873-2238}}

\affil[1]{Lunar and Planetary Laboratory, University of Arizona, Tucson, AZ, USA, https://ror.org/03m2x1q45}

\setcorrespauth{Michael S. Phillips}
\setcorrespmail{phillipsm@arizona.edu}

% ------ BIB file -----------------------------------------------------
\addbibresource{references_final_rumil.bib}
% Rumil bib carries ISSN, full dates, and complete author lists; print as in v5 (19 names + et al. above 20).
\ExecuteBibliographyOptions{maxbibnames=20,minbibnames=19}
\AtEveryBibitem{\clearfield{issn}\clearfield{month}\clearfield{day}\iffieldundef{doi}{}{\clearfield{url}}}

% ------ Date ---------------------------------------------------------
\date{}

% ------ Document -----------------------------------------------------
\begin{document}
\maketitle
\thispagestyle{firstpage}

\correspondingauthor

\submissiondate{}

\mstype{\articletype}

\summarycharactercount{(345 characters)}

\IfStrEq{\articletype}{LETTER}{
\letterwordcount{($\le$ 4500)}
\methodswordcount{($\le$ 3000 words)
}}{
\abstractwordcount{(250 words)}
}

\stitle{\insertshorttitle}

% ------ Summary ------------------------------------------------------
\summary{I propose a model for hollow formation in which impact heat drives sulfur from sulfides into cooler regolith, and the charged regolith collapses as it loses mass. A reservoir of sulfur remains stable below hollows today and sublimates in sunlight, causing recent and ongoing activity. A global statistical analysis of hollows supports the model.}

% ------ Abstract -----------------------------------------
\begin{abstract}
Hollows are among the most enigmatic features in the solar system. They are bright, flat-floored depressions on Mercury, tens of meters deep, and among the youngest landforms on its surface. They resemble sublimation pits or karst features, but form in the apparently rocky regolith of the Sun's nearest neighbor. The volatile and the process that form hollows are unknown. Here I analyze a global hollow catalog to understand what factors control hollow formation. Hollows form where impacts expose older crust beneath the volcanic plains and are more common in fresh craters and at the hot longitudes near the equator. Craters inside older craters and craters with explosive volcanic vents host hollows more often than other craters. Hollow depth increases with hollow size but varies little with surface temperature or host-crater size. I propose a model in which impact heat decomposes sulfides derived from the older crust or upper mantle, and sulfur vapor charges cooler regolith with native sulfur. While the crater is hot, sulfur vapor escapes from the lower part of this charged layer, and the porous, weakened regolith compacts under the weight of the overlying rock, analogous to karst and thermokarst. Complete compaction produces the measured depths over a plausible range of temperature and bulk porosity values. Native sulfur remains quasi-stable long-term beneath a refractory lag in the cooler upper part of the layer. Where small (micrometeoroid) impacts, mass wasting, or scarp retreat expose the volatile-infused regolith, daytime sublimation keeps hollows fresh and accounts for recent or ongoing activity.
\end{abstract}

% ------ License ------------------------------------------------------
\license{}

\section{Introduction}
Hollows are shallow, irregular, rimless depressions on Mercury with flat floors, bright interiors, and bright halos \parencite{blewett2011,blewett2016,blewett2018}. They are tens of meters deep, lack superposed craters, and are among the youngest landforms on the planet, which indicates recent or ongoing loss of a volatile phase \parencite{blewett2011,wang2020}; one hollow on the rim of Lermontov crater may have grown during the MESSENGER mission \parencite{speyerer2022}. Hollows occur mostly in impact craters, on central peaks, peak rings, walls, floors, and rims, commonly in low-reflectance material excavated from depth \parencite{blewett2013,thomas2014,thomas2016,xiao2013,bickel2025}. The colors and spectra of their bright floors and halos are consistent with calcium and magnesium sulfides, which has led some to propose that sulfides are the substance lost to form hollows \parencite{vilas2016,thomas2016,barraud2020,emran2025}. The identification is contested: chlorides fit the spectra of hollows in Tyagaraja crater better than sulfides do \parencite{lucchetti2021}, and the dark spots that surround young hollows did not change over 30 months, which is evidence against unstable MgS and CaS as their main component \parencite{wang2022}.

A model of hollow formation must explain several, and some apparently contradictory, observations at once (Table~\ref{tab:constraints}). Hollows form in material that impacts excavate from beneath the young volcanic plains \parencite{blewett2013,thomas2014}. Hollow floors are free of craters, yet hollows occur in craters hundreds of millions of years old \parencite{blewett2016,wang2020}. Hollows are more common at sunlit sites, including the hot longitudes at 0$^\circ$ and 180$^\circ$E, which face the Sun at perihelion \parencite{detoffoli2024,deutsch2025}, yet hollows have similar depths, tens of meters, across the planet \parencite{vaughan2012,blewett2013,thomas2014,blewett2016}. Proposed models for hollow formation differ in the source and identity of the volatile and in what drives its loss. Impacts may excavate volatile-rich material that sublimates until a lag covers the floor \parencite{blewett2013,blewett2016} or until deepening stops at a silicate residue \parencite{emran2025}. Impact heat may drive sulfur-rich gas toward the surface \parencite{phillips2021}, where it may react with surface minerals \parencite{renggli2022}, and sulfur may migrate between the exosphere and the subsurface \parencite{verkercke2025}. A chloride-rich layer may form as impact melt differentiates \parencite{vaughan2012,lucchetti2021}. A buried volatile-rich layer, possibly halite, may flow and behave like an ice at Mercury environmental conditions \parencite{rodriguez2020,rodriguez2023,domingue2025}. Micrometeoroid bombardment and insolation may drive the loss \parencite{bickel2025}, and volatile loss may drive long-lived mass wasting \parencite{wright2020}. Each of these models accounts for part of the observations listed in Table~\ref{tab:constraints}.

Here I present a complete model of hollow formation, from the source of the volatile to the degradation of the hollow, and test each step against available observations. The model is a revision of that of \textcite{phillips2021} and includes the sulfidation proposed by \textcite{renggli2022}. For the tests, I used the global hollow catalog of \textcite{bickel2025}, normalized by imaging coverage and detection probability, together with crater and terrain data, the explosive volcanic vents of \textcite{pegg2021}, thermal models, the shadow-length depths of \textcite{blewett2016}, and shadow-length depths that I measured.

\begin{table}[!ht]
    \mytablestyle
    \caption{{\bf Observations that constrain hollow formation.} The last column gives the step of the model (Section~\ref{sec:model}) that explains each observation.}
    \scriptsize{
        \begin{tabularx}{\textwidth}{>{\raggedright\arraybackslash}p{4.0cm} >{\raggedright\arraybackslash}X >{\raggedright\arraybackslash}p{5.0cm}}
            \toprule
            \textbf{Observation} & \textbf{Evidence} & \textbf{Step} \\
            \hline
            Hollows form in older crust excavated from beneath young volcanic plains & Scarce on the plains; in Caloris, hollow fields only in craters whose excavation reaches below the plains; densest on central peaks and peak rings (Section~\ref{sec:source}; \textcite{blewett2013,thomas2014}) & Source; charging \\
            Craters inside older craters host hollows more often, and craters inside hollow-bearing craters most often & Section~\ref{sec:nesting}; \textcite{giroud2024,detoffoli2024} & Charging and activity: a later impact reheats and exposes the parent's charged layer \\
            Hollow spectra are consistent with Ca and Mg sulfides and graphite & \textcite{vilas2016,helbert2013,thomas2016,lucchetti2018,barraud2020,emran2025}; contrary result in \textcite{lucchetti2021} & Source; sulfidation during charging \\
            Hollows are active in host craters more than 10$^8$~yr old & Hollow floors are free of craters for $\leq$ $\sim$10$^5$~yr \parencite{wang2020} in hosts up to $>$ 1.7~Gyr old \parencite{blewett2013,banks2017} & Storage under a refractory lag \\
            Fresh craters host hollows more often, and hollows of a given size are as deep in fresh as in degraded craters & Sections~\ref{sec:age} and \ref{sec:depth}; \textcite{detoffoli2024} & Dormancy: exposed layers are spent or buried; collapse during charging sets depth \\
            Craters at the hot longitudes host hollows more often near the equator and to 30$^\circ$N, and hollows are more common on equator-facing walls at high northern latitudes & Section~\ref{sec:sun}; \textcite{detoffoli2024,deutsch2025} & Activity, modulated by sunlight \\
            Hollow floors are flat, and hollows in one crater can share a common maximum depth & \textcite{vaughan2012,blewett2013,thomas2014,blewett2016} & Charging and collapse: the charged layer collapses beneath the floor \\
            Depth increases with hollow size, is set locally, and varies little with temperature and host size & Section~\ref{sec:depth} & Charging and collapse: vapor pressure and gravity set the base of the layer, and the volume lost on collapse sets the depth, if base temperatures are similar among craters; $\phi$ may differ between hollow fields (Section~\ref{sec:depthcontrol}) \\
            Hollows have bright halos & \textcite{blewett2011,blewett2013,blewett2016} & Activity: sulfur vapor reacts with the surrounding regolith \\
            Hollows have irregular outlines and occur in fields & \textcite{blewett2011} & Charging: vapor spreads laterally through fractured regolith; fields form as hollows widen by scarp retreat and merge \\
            Craters with explosive volcanic vents host hollows more often & Section~\ref{sec:vents}; \textcite{pegg2021} & Charging: magmatic intrusions supply heat or sulfur \\
            \hline
        \end{tabularx}}
    \label{tab:constraints}
\end{table}

\FloatBarrier
\section{Data and Methods}
\label{sec:methods}

\subsection{Hollow Catalog and Imaging Coverage}
\label{sec:catalog}
\textcite{bickel2025} released 19,110 hollow detections, found by a deep-learning detector in 4,915 images from the narrow-angle camera of the MESSENGER Mercury Dual Imaging System \parencite[MDIS;][]{hawkins2007}. I recomputed each detection's position with the MDIS narrow-angle camera model in the Integrated Software for Imagers and Spectrometers \parencite[ISIS;][]{isis2023} and reconstructed MESSENGER navigation data \parencite{spice_messenger}. For hollows detected in two images at $<$ 60~m/pixel, the median distance between the two detections falls from 1.90~km for the released positions to 0.41~km (Section~S1; Figure~S1). I kept detections with confidence $\geq$ 0.75 in images with solar incidence $<$ 85$^\circ$ and merged repeat detections into 6,731 hollows. Many large detections outline crater bowls or sunlit walls, so I screened 768 detections on image crops, chosen among detections at least a third as wide as the crater they lie in, large detections in fields of at most three hollows, detections 1.0--1.5~km across in host craters with at most two hollows, and a random control sample, and I excluded the 368 that are crater bowls, walls, or other features that are not hollows (Section~S1). The remaining 6,363 hollows form 973 hollow fields, groups of hollows within 20~km of one another.

Hollow counts depend on where MESSENGER imaged and on how well the detector found hollows in each image. I modeled the probability that the detector finds a hollow as a function of the hollow's diameter in pixels and the solar incidence angle, using 85,602 hollow--image pairs, and combined it with the footprints of all images to compute the number of hollows or fields expected in any region if hollows were distributed uniformly (Section~S2; Figure~S2). The same model gives the ``completeness'' of each map cell, the probability that a hollow 1~km across is detected in at least one image of the cell. At the sites of cataloged hollows, 1~km is about the size at which this probability reaches one half (median 54\%; Section~\ref{sec:sfd}). The results change little for reference sizes from 0.3 to 1~km, and the hosting results hold with completeness averaged over the size distribution of hollows in place of a single size (Section~S2). I compared regions with different imaging coverage by the ratio of observed to expected counts.

The same detection model gives an estimate of how many hollows of each size exist at the sites MESSENGER imaged, including those the detector missed. For each cataloged hollow, the model gives the probability that a hollow of any diameter at that site would be detected in at least one of the images that cover it, using each image's own pixel scale and solar incidence. I fitted power-law size distributions over ranges of hollow diameter by maximum likelihood, conditional on the locations of the cataloged hollows, so that each hollow's diameter is weighed against the range of sizes the images at its site could detect. I also counted hollows corrected for detection by weighting each cataloged hollow by the inverse of its probability of detection, with 95\% intervals from 300 bootstrap resamples of hollows (Section~S2).

\subsection{Craters, Terrain, and Temperature}
\label{sec:craters}
I assigned each hollow to the smallest crater $\geq$ 5~km in the catalog of \textcite{herrick2018} whose 1.1 radii contain it. A crater hosts hollows when at least one hollow is assigned to it, which gives 534 host craters among 31,571 craters. In the 114 host craters 30--250~km across with at least five hollows, I traced the outlines of peaks and the foot of the walls by hand on the 166~m/pixel MDIS mosaic \parencite{denevi2018} and the 665~m/pixel global topography \parencite{becker2016}, and counted material slumped from the walls as wall (Figure~S3). In 3 of these craters the peak could not be distinguished from the floor, so the terrain tests use the other 111, which hold 3,323 of the 3,761 hollows (88\%) assigned to the 321 host craters 30--250~km across. I divided the interior of each of these craters into peak, floor, and wall terrain and weighted the area of each terrain by the completeness of the map cells it covers to account for area that MESSENGER imaged poorly (Section~S4). I mapped the northern smooth plains \parencite[Section~S6;][]{head2011} from the lineaments of a global tectonic map \parencite{klimczak2025}, and I estimated the excavation depth of craters inside the Caloris basin from crater scaling \parencite[Section~S7;][]{croft1985,melosh1989}. I took the positions of 309 candidate explosive volcanic vents from the catalog of \textcite{pegg2021} and assigned each vent to the smallest crater $\geq$ 5~km whose rim contains it (Section~S15). For craters $\geq$ 40~km, I also used the degradation classes of \textcite{kinczykMorphologicalEvaluationCrater2020}, matched to the catalog of \textcite{herrick2018} by position and diameter (Section~S11).

As in \textcite{vasavada1999}, I computed surface temperature with a one-dimensional thermal model that includes Mercury's 3:2 spin--orbit resonance and eccentric orbit; unlike theirs, my model uses constant thermal conductivity (Section~S5). Because of the resonance, noon at the hot longitudes, 0$^\circ$ and 180$^\circ$E, falls at perihelion, and noon at the warm longitudes, 90$^\circ$ and 270$^\circ$E, falls at aphelion. I define $\Delta T_\mathrm{w}$ as the modeled mean annual surface temperature at a site minus that at the warm longitude of the same latitude; $\Delta T_\mathrm{w}$ reaches 62~K at the hot longitudes on the equator. To test for an excess of hollows at the hot longitudes, I fit logistic regressions of whether a crater hosts hollows on $\Delta T_\mathrm{w}$, log$_{10}$ crater diameter, their interaction, and absolute latitude. Holding latitude fixed removes most of the effect of imaging coverage, which varies more with latitude than with longitude. Coverage still varies with longitude (Section~S3), so I also refit with completeness at the crater center as a control. Where I write that craters host hollows some number of times as often as other craters, that number is an odds ratio from such a fit, unless stated otherwise. Uncertainties are 95\% intervals; intervals on odds ratios are Wald intervals from the logistic fits, and other intervals come from 2,000 bootstrap resamples.

\textbf{Crater containment graph.} To test whether craters that form inside older craters, and so expose material that the older impact already raised, host hollows more often, I linked Mercury's craters by containment into a graph of stratigraphic relationships (Section~S14). The nodes are the 31,571 craters $\geq$ 5~km of \textcite{herrick2018} and the 52 basins of \textcite{orgel2020} that are rated certain or probable and do not duplicate a cataloged crater. A crater whose center lies inside the rim of a crater or basin at least twice its diameter, so that the smaller crater lies mostly inside the larger one, formed after that feature. I did not order craters of similar size whose rims overlap, because the catalog of \textcite{herrick2018} does not record which rim is better preserved where they cross. Each crater's parent is the smallest feature that contains it, so the craters form a directed graph in which each crater has at most one parent and every edge runs from an older to a younger crater. I dropped from the tests the inner members of 17 nested pairs that may be peak rings, which leaves 31,554 craters. I assigned each crater to its parent's peak ($\rho <$ 0.2), floor (0.2 $\leq \rho <$ 0.6), or wall ($\rho \geq$ 0.6) by the distance $\rho$ of its center from the parent's center, in parent radii. The graph, the code, and an interactive map of the graph are archived at Zenodo \parencite{zenodo_cratergraph}.

For each crater I estimated a first-order depth of origin, the deepest pre-impact level it exposes, from its excavation depth or stratigraphic uplift \parencite{croft1985,melosh1989,cintala1998}, the depth of origin of its parent, and the thickness of fill on the parent's floor (Section~S14). I repeated every depth-of-origin fit with the terrestrial and the lunar uplift relations of \textcite{cintala1998}. Both relations were fitted to craters with central peaks, so beneath peak-ring craters and basins I set the uplift near the maximum depth of excavation, the depth from which the rock of lunar peak rings came \parencite{baker2015,kring2016}.

In these tests a crater counts as a host when at least one hollow assigned to it lies within 0.9 of its radius, so hollows on a parent's floor just outside a smaller crater do not count for that crater; 446 craters qualify. I fit logistic regressions of hosting on nesting, position in the parent, depth of origin, number of ancestors, whether the parent hosts hollows, and whether the crater contains a volcanic vent. Every fit includes log$_{10}$ diameter, absolute latitude, completeness at the crater center, a northern-plains flag, and $\Delta T_\mathrm{w}$ as controls, and the fits for odds ratios also include the logarithm of the fraction of craters within 500~km that host hollows. For the comparisons by the Akaike information criterion (AIC), I used the 532 host craters of the assignment above and omitted the regional hosting fraction. The AIC measures how well a model fits the data, with a penalty for model complexity of 2 for each fitted parameter. A lower AIC indicates a better model, and a model whose AIC exceeds that of the best model by more than 10 has essentially no support relative to it \parencite{burnham2004}. I also fit a conditional logistic regression to craters inside the same parent, and I compared nested craters with isolated craters of similar size and completeness within 300~km (Section~S14).

\subsection{Hollow Depths}
\label{sec:depthmethods}
I used two sets of shadow-length depths. \textcite{blewett2016} measured 2,518 hollow depths in 565 images; their depth table lists 527 distinct images. I measured depths automatically in the raw MDIS narrow-angle images (Section~S9), from the length of the shadow that the sunward wall casts along the Sun's direction inside each detection of \textcite{bickel2025}, corrected for foreshortening and divided by the tangent of the local solar incidence angle $i$. I used the finest detection of each hollow in images at $\leq$ 20~m/pixel, 58--81$^\circ$ incidence, and $\leq$ 50$^\circ$ emission, which gives 1,558 depths in 592 images. I also measured by hand the shadows of 135 hollows 40~m to 1~km across in images at $\leq$ 7~m/pixel, to validate the automatic depths and to test how depth changes with hollow size. To compare my measurements with theirs, I measured by hand the shadows at 178 of the points that \textcite{blewett2016} measured, without seeing their depths (Section~S9). The shortest measurable shadow corresponds to a minimum measurable depth that scales with pixel scale divided by $\tan i$. I therefore regressed the natural logarithm of depth on inverse mean annual temperature with ln pixel scale and ln tan $i$ as controls and standard errors clustered by image (Section~S9), and I repeated each regression on my depths with the logarithm of hollow diameter as a further control. I compare the result with a sublimation front limited by vapor diffusion through a porous lag, using elemental sulfur as a proxy for a volatile with hollow-forming properties \parencite{phillips2021}. Such a front deepens with the square root of time and of the saturation vapor density \parencite{schorghofer2018}, so, neglecting the weak $1/T$ factor in the vapor density, the logarithm of its depth falls with inverse temperature at a slope of $-$4,514~K, half that of a fit to the tabulated vapor pressures of sulfur \parencite{lide2007}. Where I compared depths between hollows or between host craters, I corrected each depth for pixel scale and solar incidence to the medians of the measurements, 14.2~m/pixel and 67.6$^\circ$ (Sections~S9 and S10). I also measured how depth correlates between pairs of hollows as a function of their separation (Section~S10), and I regressed depth on host diameter, on host degradation state (Section~S11), on MESSENGER X-ray elemental ratios \parencite[Section~S9;][]{nittler2020}, and on color (Section~S9).

\FloatBarrier
\section{Results}
\label{sec:results}

\subsection{Small Hollows Are Either Missing from the Catalog or Rare, and Large Hollows Are Rarer Than a Power Law Predicts}
\label{sec:sfd}
After correction for the hollows that the detector of \textcite{bickel2025} missed, the size distribution of hollows follows a power law between 0.3 and 3~km and steepens above 3~km (Figure~\ref{fig:sfd}a). Between 0.3 and 3~km, the number of hollows larger than $D$ is proportional to $D^{-1.4}$ to $D^{-1.6}$. Above 3~km, the exponent is 3.3 (2.9--3.6), a rollover that could reflect an upper limit on hollow size or a limitation of the detector. After correction for the hollows the detector missed, the detection model predicts 105 (84--128) hollows larger than 5~km, against 360--490 if the power law between 0.3 and 3~km continued. At these sizes a median of 84--86\% of hollows would be detected (Figure~\ref{fig:sfd}b), so the correction is small, and the shortage of large hollows does not come from missed detections or from detections that outline only part of a hollow (Section~S2). As \textcite{bickel2025} did, I excluded the 64 detections larger than 10~km, which outline whole host craters, in addition to the screened detections (Section~\ref{sec:catalog}).

At smaller sizes the catalog appears far from complete. At the sites of cataloged hollows, the median probability that a hollow is detected in at least one image is 54\% for a hollow 1~km across and 4\% for one 0.3~km across (Figure~\ref{fig:sfd}b). The detector misses small hollows for two reasons. First, a hollow 0.3~km across spans fewer than 10 pixels in 79\% of the images that cover cataloged hollows (median pixel scale 58~m), and the detector reaches half its highest detection rate only for hollows about 35 pixels across (Figure~S2b). Second, the detector misses many hollows that other images resolve: a cataloged hollow is found in only 5\% of the other images that cover it, and in 15--17\% of those in which it spans 60--200 pixels. The catalog holds 5,676 hollows larger than 0.3~km. Weighting each by the inverse of its probability of detection gives 21,300 (20,400--22,100) hollows larger than 0.3~km. This is a lower limit, because sites where no image could detect a hollow carry no weight. Extending the power law from 1~km down to 0.3~km gives 25,000--32,000 hollows larger than 0.3~km, so the catalog holds about one in four to one in six of all hollows $>$ 0.3~km on the surface the searched images cover. From the recall of their detector, \textcite{bickel2025} estimated about 34,000 hollows in MESSENGER's narrow-angle images after removing repeats in overlapping images, consistent with this estimate. At the sites imaged finely enough to detect a 0.3~km hollow with a probability of at least 30\%, counts keep rising toward smaller sizes, with an exponent of 2.3 (1.8--2.9) between 0.3 and 1~km (149 hollows), so I found no sign of a minimum hollow size down to 0.3~km.

\begin{figure}[!ht]
    \centering
    \includegraphics[width=15cm]{figures/fig_sfd_v3_v3cat.png}
    \caption{{\bf Size distribution of hollows.} (a) Number of hollows larger than diameter $D$: cataloged (orange), and corrected for missed hollows by weighting each cataloged hollow by the inverse of its probability of detection (blue; 95\% band from 300 bootstrap resamples of hollows), which is a lower limit at small $D$. Lines are power laws with exponents of 1.4--1.6, the range fitted over 0.3--3~km and drawn through the corrected count at 1~km (dashed), and 3.3, fitted above 3~km (dotted). The shaded region between the cataloged and corrected counts estimates the hollows the detector missed. Hatching marks the hollows missing below 1~km if the power law continues to small sizes, and the shortfall from the power law above 3~km, which may reflect an upper limit on hollow size. (b) Probability that a hollow of diameter $D$ at the site of a cataloged hollow is detected in at least one image: median and interquartile range over the sites.}
    \label{fig:sfd}
\end{figure}

\subsection{Hollows Form Where Impacts Expose the Crust Beneath the Plains}
\label{sec:source}
Hollows are rare on the northern smooth plains \parencite{blewett2013,thomas2014}, and the scarcity holds with imaging coverage controlled. In cells imaged at $\leq$ 40~m/pixel, hollow fields inside the plains are 0.35 times as dense as outside at 40--60$^\circ$N and 0.13 times as dense at 60--80$^\circ$N (Figure~\ref{fig:exposure}a, b; Section~S6).

\begin{figure}[!ht]
    \centering
    \includegraphics[width=15cm]{figures/fig1_exposure_v10.png}
    \caption{{\bf Hollows form where impacts expose older crust.} (a) North polar equal-area map from 30$^\circ$N, with 0$^\circ$E at the bottom, on the 166~m/pixel MDIS mosaic \parencite{denevi2018}, showing the finest pixel scale of the images that cover each area (up to 150~m/pixel), the smooth-plains mask, the part of the Caloris basin rim north of 30$^\circ$N \parencite[dashed;][]{orgel2020}, and hollow fields scaled by number of hollows. (b) Hollow fields per 10$^6$~km$^2$ imaged at $\leq$ 40~m/pixel, outside the plains with all images, with only images at 60--75$^\circ$ incidence, and without the province at 90--150$^\circ$E, 40--70$^\circ$N, and inside the plains; 95\% Poisson intervals. (c) Fraction of craters hosting $\geq$ 3 hollows against crater diameter, for craters inside Caloris and isolated craters, both with completeness at the crater center $\geq$ 0.2, with 95\% binomial intervals and Caloris counts. The band marks diameters at which excavation reaches 2.5--3.5~km, the plains thickness of \textcite{ernst2015}. (d) Hollow density on peaks, floors, and walls of the 111 host craters 30--250~km across whose terrain I traced by hand, by crater setting, with area weighted by completeness; 95\% bootstrap intervals over craters (Table~S2).}
    \label{fig:exposure}
\end{figure}

Inside the Caloris basin, the volcanic plains are 2.5--3.5~km thick \parencite{ernst2015}, and the impacts that formed larger craters excavated darker material from beneath them \parencite{murchie2008}; \textcite{blewett2013} reported that hollows in Caloris lie in this excavated material. No crater inside Caloris whose excavation falls short of 2.5~km hosts three or more hollows (0 of 144), and 14 of 35 craters whose excavation reaches deeper do (Fisher $p$ = 9.8 $\times$ 10$^{-12}$; Figure~\ref{fig:exposure}c). Small craters elsewhere on Mercury also rarely host three or more hollows (2 of 840 isolated craters 10--30~km across with completeness $\geq$ 0.2), so the absence of hollow fields in the shallow Caloris craters is expected from their size alone, and the effect of the plains on the preponderance of hollows can be tested only among larger craters. Among craters 30--35~km across, 3 of 5 Caloris craters host three or more hollows, compared with 1 of 58 isolated craters of similar completeness elsewhere ($p$ = 0.001; Section~S7). Material that only craters excavating below the plains reach therefore appears to be required for hollow fields in Caloris, and the low-reflectance material beneath the plains is a candidate for this material.

Within host craters, hollows are densest where impacts raise the deepest material. With each unit of area weighted by its completeness, central peaks and peak rings have 7.2 (4.3--12.2) times the hollow density of walls in the 111 traced host craters (Figure~\ref{fig:exposure}d; Section~S4). Floors have 1.6 (0.98--2.7) times the hollow density of walls, a ratio whose interval includes 1. \textcite{detoffoli2024} likewise found 34\% of hollow-bearing crater locations on central peaks.

Within peaks, hollows are least dense near the base. Peak terrain within 2~km of the floor has 0.48 (0.39--0.59) times the hollow density of the rest of the peak, where each crater's hollows are compared with that crater's peak area weighted by completeness (1,448 hollows on the peaks of 92 craters; Section~S4). Above that level, hollow density does not change with height (Section~S4).

Hollows also occur outside craters on the knobs that surround Caloris \parencite{wright2020}. Of the 2,011 knobs that \textcite{wright2020} mapped, 19 have 24 cataloged hollows in 15 fields. With area weighted by completeness, the knobs have 4.0 (1.9--7.8) times the hollow density of the terrain 1.5--3 knob radii from them (Section~S17). The knobs are as old as Caloris, and degraded craters (degradation state is a proxy for crater age) host hollows less often (Section~\ref{sec:age}), so I also compared the knobs with the peaks of the most degraded craters, Classes~1--2 of \textcite{kinczykMorphologicalEvaluationCrater2020}. I traced the peaks by hand in all 108 of these craters that host hollows and in a random sample of 60 of the 2,052 that do not (Section~S17). The knobs have about the same hollow density as these peaks (ratio 0.8; 0.3--1.6).

\FloatBarrier
\subsection{Craters Inside Older Craters Are More Likely to Host Hollows}
\label{sec:nesting}
Of the 31,554 craters in the containment tests, 14,916 (47\%) formed inside a crater or basin at least twice their diameter: 9,788 on the parent's wall, 4,565 on its floor, and 563 on its peak. \textcite{giroud2024} found that hollows in the most degraded host craters lie mostly in small superposed craters, and \textcite{detoffoli2024} found hollows in older craters with deep-seated features; I tested these observations across all cataloged craters with the containment graph.

Craters that formed inside an older crater or basin host hollows 1.68 (1.35--2.09) times as often as isolated craters (Figure~\ref{fig:nesting}). The ratio is 1.42--2.05 in every subset that includes all parents, including craters with at least three hollows, the 1,252 craters with completeness $\geq$ 0.5, and craters with degradation state controlled, and the lower bound of its interval exceeds 1 in each (Table~S17). Nested craters host hollows 1.7 times as often as isolated craters of similar size and completeness within 300~km (0.83\% against 0.48\%; Section~S14).

\begin{figure}[!ht]
    \centering
    \includegraphics[width=15cm]{figures/fig_nested_examples_v3cat.png}
    \caption{{\bf Craters inside two hollow-bearing craters.} Praxiteles (183~km; 27.2$^\circ$N, 299.6$^\circ$E) and Lermontov (162~km; 15.3$^\circ$N, 311.0$^\circ$E) on the 166~m/pixel MDIS global mosaic \parencite{denevi2018}. Dashed white circles mark the parent rims, and arrows run from each parent's center to the craters $\geq$ 5~km that formed inside it. Gold circles mark craters with at least one hollow inside 0.9 radii, blue circles craters without, and yellow dots the georeferenced hollows. Four of 6 craters inside Praxiteles and 3 of 5 craters inside Lermontov host hollows, compared with 1.4\% of all craters $\geq$ 5~km. The hollows of both craters have been studied in detail \parencite{pajola2021,speyerer2022,galiano2026}.}
    \label{fig:nesting}
\end{figure}

Craters inside a parent that hosts hollows host hollows 2.98 (1.98--4.50) times as often as craters inside a parent without hollows. Here a parent hosts hollows when at least one hollow assigned to the parent itself, and not to a smaller crater inside it, lies within 0.9 of its radius. The ratio is 2.28 (1.25--4.14; $p$ = 0.007) with degradation state controlled (1,295 such craters $\geq$ 10~km with a degradation state), and 2.62 (1.73--3.97) when a parent counts as a host if any hollow is assigned to it. By that definition, 97 of 2,579 craters (3.8\%) inside hollow-bearing parents host hollows, compared with 40 of 7,513 (0.53\%) inside other parents (Section~S14). MESSENGER imaged known hollow sites at high resolution, but the inheritance is not an effect of better imaging: among the 414 craters inside parent craters with completeness $\geq$ 0.5, craters inside hollow-bearing parents host hollows 4.53 (1.71--11.9) times as often, although only 39 of them host hollows.

Position within the parent and depth of origin add little to nesting. Within the same parent, craters on the peak or floor host hollows 1.25 (0.84--1.86; $p$ = 0.27) times as often as craters on the wall, a ratio whose interval includes 1. In fits without the regional hosting fraction, nesting improves the model fit (AIC 38 lower). Added to the logarithm of a crater's own excavation or uplift depth, the logarithm of its depth of origin lowers the AIC by 22 ($p$ = 4 $\times$ 10$^{-7}$) with either uplift relation, but every depth-of-origin model has an AIC at least 25 higher than the model with nesting alone, so no depth-of-origin model has support relative to a model with nesting only (Section~S14). Nesting therefore predicts hosting far better than the depth a crater reaches by forming inside another crater, although depth of origin adds a small, significant effect.

\FloatBarrier
\subsection{Craters with Volcanic Vents Host Hollows More Often}
\label{sec:vents}
Of the 309 candidate explosive vents in the catalog of \textcite{pegg2021}, 243 lie within the rims of cataloged craters $\geq$ 5~km, and 148 craters in the containment tests contain at least one vent. Of these 148 craters, 34 (23\%) host hollows, compared with 1.4\% of all craters $\geq$ 5~km. With the controls of the nesting tests, craters with vents host hollows 3.46 (1.95--6.13) times as often as craters without vents, isolated craters with vents 2.51 (1.19--5.29) times as often as other isolated craters, and craters $\geq$ 20~km with vents 3.31 (1.86--5.89) times as often as other craters of that size (Section~S15). Hollows in the 13 host craters with vents that have depths in the data of \textcite{blewett2016} (405 hollows) are 0.96--1.28 times as deep as hollows in other host craters, with pixel scale, incidence, and latitude controlled, and in my depths they are 1.14 (0.95--1.37) times as deep (Section~S15).

\subsection{Hollows Are More Common Where Insolation Is Higher}
\label{sec:sun}
Hollows are more common near the hot longitudes \parencite{detoffoli2024}, and \textcite{thomas2014} found that the hollowed fraction of the surface peaks near the hot longitude at 0$^\circ$E, with no comparable increase at 180$^\circ$E. With imaging coverage and detection probability controlled, and with fields near the Caloris basin excluded, hollow fields within 15$^\circ$ of the equator are 1.64 (1.26--2.14) times as common, relative to expectation, near the hot longitudes as near the warm longitudes (Figure~\ref{fig:insolation}a; Table~S5). This excess does not come from the crater population. Craters $\geq$ 10~km are 16\% less dense near the hot longitudes than near the warm longitudes, and within 15$^\circ$ of the equator, craters near the hot longitudes host hollows 2.2 (1.5--3.2) times as often as craters of the same size, latitude, and completeness near the warm longitudes (Section~S8). At 15--30$^\circ$N, craters near the hot longitudes also host hollows more often, by a factor of 2.6 (1.5--4.4); at 30--45$^\circ$N, the factor of 1.5 (0.9--2.4) has an interval that includes 1.

Counted as hollow fields or as individual hollows, the pattern reverses away from the equator. At 45--60$^\circ$N, hollow fields are 0.39 (0.24--0.57) times as common at the hot longitudes as at the warm longitudes, and about 40\% of this reversal, on a logarithmic scale, comes from the hollow-rich province at 90--150$^\circ$E (Sections~\ref{sec:joint} and S8). At 15--75$^\circ$N, 2,373 hollows lie near the warm longitudes and 1,096 near the hot longitudes. Eight of the ten host craters with the most hollows, which together contain 50\% of the hollows at these latitudes, lie near the warm longitudes, and without these ten craters, 953 hollows lie near the hot longitudes and 780 near the warm longitudes (Section~S8). These few hollow-rich craters account for the high hollow densities near the warm longitudes found by \textcite{bickel2025}. At high northern latitudes, hollows are also more common on equator-facing walls (Section~S8), as \textcite{bickel2025} found for equator-facing slopes globally, although \textcite{thomas2014} found a dominant slope aspect in only 8\% of groups of hollows.

The odds that a crater hosts hollows rise with $\Delta T_\mathrm{w}$. Within 60$^\circ$ of the equator, with diameter, latitude, and completeness controlled, they rise by a factor of 1.37 (1.25--1.50) per 10~K of $\Delta T_\mathrm{w}$ (Section~S8). The rise appears steeper in smaller craters (Figure~\ref{fig:insolation}b; Section~S3), but the change with crater size is not statistically significant: the interaction between $\Delta T_\mathrm{w}$ and crater size has $p$ = 0.061, or 0.21 with completeness at the crater center as a control.

The freshness of hollows depends on insolation more abruptly than their occurrence does. \textcite{deutsch2025} classified hollows as bright and sharp (Stage 1), muted (Stage 2), or faded (Stage 3). Among hollows classified in images with solar incidence below 70$^\circ$, Stage 1 hollows make up 56--71\% of the hollows where $\Delta T_\mathrm{w}$ is below 45~K and 96--100\% where it is above 50~K. In a logistic fit with solar incidence, pixel scale, emission angle, and latitude as controls, the odds that a hollow is classified as Stage 1 rise 3.0-fold (1.5--6.2) in a step at $\Delta T_\mathrm{w}$ = 44~K. A model with this step has an AIC 27 lower than one in which the log-odds rise linearly with $\Delta T_\mathrm{w}$ (Section~S8). Classifications from images at higher incidence depend on the incidence and show no such step.

\begin{figure}[!ht]
    \centering
    \includegraphics[width=15cm]{figures/fig3_insolation_v9_v3r.png}
    \caption{{\bf Hollows are more common where insolation is higher.} (a) Ratio of observed to expected hollow fields in 30$^\circ$ longitude bins within 15$^\circ$ of the equator (217 fields; fields within 1,200~km of the Caloris basin center excluded). The line and band show the fit $0.99 + 0.50 \cos 2\lambda$, where $\lambda$ is east longitude, and the 95\% interval on its amplitude (0.34--0.65). (b) Odds ratio of hosting hollows per 10~K of $\Delta T_\mathrm{w}$ against crater diameter, for 27,271 craters $\geq$ 5~km within 60$^\circ$ of the equator (485 hosts). The line and band come from a logistic fit with a $\Delta T_\mathrm{w}$ $\times$ log diameter interaction, which is not significant ($p$ = 0.061), and the points from fits in six size bins, labeled with host counts.}
    \label{fig:insolation}
\end{figure}

\FloatBarrier
\subsection{Fresh Craters Host Hollows More Often}
\label{sec:age}
\textcite{herrick2018} classified unfilled craters $>$ 10~km as fresh, standard, or degraded by the sharpness of their rims, the preservation of their interior structure and ejecta, and the presence of rays. Among the 4,216 craters $\geq$ 10~km with a degradation state in this catalog, fresh craters have 4.9 (3.3--7.2) times the odds of degraded craters of hosting hollows, and standard craters have 1.9 (1.4--2.6) times the odds, at fixed diameter and latitude (Section~S11). \textcite{detoffoli2024} reported qualitatively that hollows are common in young craters and rare in old, degraded ones.

\textcite{kinczykMorphologicalEvaluationCrater2020} classified all craters $\geq$ 40~km, including the filled craters that \textcite{herrick2018} left unrated, on a five-point scale from Class~1, with discontinuous rims, no ejecta, and no central peak, to Class~5, with crisp rims and bright rays. At fixed diameter and latitude, the odds of hosting rise 2.9 (2.5--3.4) times per class (Section~S11).

Host craters range from $<$ 0.3~Gyr \parencite{blewett2016,banks2017} to Calorian basins older than 1.7~Gyr \parencite{blewett2013,banks2017}, and hollow floors are free of craters for at most $\sim$10$^5$~yr \parencite{wang2020}. Stage 1 hollows classified in images at solar incidence below 70$^\circ$ lie in 69 host craters of Classes~1--3 (Section~S11). \textcite{kinczykMorphologicalEvaluationCrater2020} linked Classes~4 and 5 to the Mansurian and Kuiperian periods, so these 69 hosts probably formed before the Mansurian period, which began about 1.7~Gyr ago \parencite{banks2017}. Conduction removes heat from a hot zone of thickness $L$ in a time of order $L^2/\kappa$. With thermal diffusivity $\kappa$ = 10$^{-6}$~m$^2$~s$^{-1}$, this time is 4 $\times$ 10$^{10}$~s, about 1,300~yr, for $L$ = 0.2~km and 2.5 $\times$ 10$^{13}$~s, about 800,000~yr, for $L$ = 5~km, and purely conductive models of craters 20--200~km across give cooling times of 10$^3$--10$^6$~yr \parencite[and references therein]{abramov2005}. At the 200~km Chicxulub crater on Earth, a hydrothermal system of water circulating through hot rock lasted 1.5 to at least 8~Myr \parencite{abramov2007,kring2020,pickersgill2026}, but Mercury's craters lack such water, so their cooling times are set by conduction. Even a crater that stayed hot for 10$^7$~yr, longer than any of these estimates, would have cooled long ago, so Stage 1 hollows in host craters older than 1.7~Gyr are ostensibly active more than 10$^9$~yr after impact heating ended.

\subsection{Hollow Depth Increases with Hollow Size, Is Set Locally, and Varies Little with Temperature, Host Size, and Composition}
\label{sec:depth}
The two sets of depths measure the same shadows alike but differ by a constant factor. At the 178 points of \textcite{blewett2016} that I measured without seeing their depths, their depths are 0.83 (0.77--0.88) times mine, by a similar factor equatorward and poleward of 60$^\circ$N (0.85 and 0.81), and the two sets rank the hollows alike (Spearman $\rho$ = 0.89). The incidence angles of the two measurements differ by a median of 0.03$^\circ$, so the difference lies in where the edges of the shadows were placed. The automatic depths agree with my hand measurements of the same hollows (median ratio 1.01; 80\% within 30\%; 114 hollows; Section~S9). The intervals below include the random scatter of individual measurements; the difference of about 20\% between observers affects absolute depths but not ratios or slopes within one set of depths (Section~S9).

Depth increases with hollow size (Figure~\ref{fig:depth}b). With pixel scale and solar incidence controlled, depth is proportional to $D^{0.57 \pm 0.07}$ in the 131 hollows I measured by hand and to $D^{0.61 \pm 0.03}$ in the 1,558 automatic depths, where $D$ is hollow diameter. Measured diameters scatter between detections of the same hollow, which flattens the fitted relation, so these exponents are lower limits; corrected for that scatter, the exponent is 0.73--0.90 (Section~S9). The median depth is 5~m for hollows 40--100~m across and 19~m for hollows 300--500~m across, so the ratio of depth to diameter falls from about 0.08 to 0.05. Neither a break in slope nor a plateau fits the hand measurements better than a single power law, so these data show no minimum hollow size and no maximum depth between 40~m and 1~km. Among hollows of one field measured in one image, the exponent is 0.37 $\pm$ 0.10 in the hand measurements and 0.49 $\pm$ 0.06 in the automatic depths, and between the means of these groups it is 0.74 $\pm$ 0.15 and 0.84 $\pm$ 0.10 (Section~S9). Depth therefore increases with size within a field, and fields of large hollows are deeper than their size alone predicts.

Depth varies little with temperature. Modeled mean annual surface temperature at the measured hollows spans 219--324~K in both sets. With pixel scale and solar incidence controlled, fitted depth changes 1.5-fold (1.2--2.1) across this range in the depths of \textcite{blewett2016} and 0.86-fold (0.66--1.12) in mine, or 0.88-fold (0.73--1.06) with hollow size also controlled, where a sulfur sublimation front driven by sunlight and limited by vapor diffusion \parencite{phillips2021} predicts a 794-fold change (Figure~\ref{fig:depth}a). The saturation vapor pressure $P$ of a volatile follows the Clausius--Clapeyron relation,
\begin{equation}
    \ln P = -\frac{\Delta H}{RT} + \frac{\Delta S}{R},
    \label{eq:cc}
\end{equation}
where $\Delta H$ and $\Delta S$ are the enthalpy and entropy of sublimation or vaporization, $R$ is the gas constant, and $T$ is temperature. To first order, the depth $h$ of a diffusion-limited front grows as $\sqrt{P}$, so between temperatures $T_1$ and $T_2$ its depth changes by a factor
\begin{equation}
    \frac{h(T_2)}{h(T_1)} = \exp\left[\frac{\Delta H}{2R}\left(\frac{1}{T_1} - \frac{1}{T_2}\right)\right].
    \label{eq:ratio}
\end{equation}
The change in depth with temperature therefore depends on the enthalpy, not on how volatile the substance is at a given temperature, which also depends on $\Delta S$. Across 219--324~K, a front driven by water ice, which is more volatile than sulfur but has a lower enthalpy, would deepen 85-fold, against 794-fold for sulfur (Table~\ref{tab:fronts}). Fitted with Equation~\ref{eq:ratio}, the measured dependences imply apparent enthalpies of 4.8 (1.6--8.1) and $-$1.8 ($-$4.7 to 1.2)~kJ~mol$^{-1}$, far below the enthalpy of any candidate volatile, so no front driven by surface temperature fits the depths, whichever of these volatiles drives it (Section~S13).

Hollows poleward of 60$^\circ$N are 0.77 (0.69--0.86) times as deep as hollows equatorward in the depths of \textcite{blewett2016} and 0.93 (0.83--1.03) times as deep in mine. Poleward hollows are smaller than equatorward hollows (median 363~m against 485~m in my depths), and 72\% of them lie in the hollow-rich province (Section~\ref{sec:sun}). With hollow size controlled, the ratio is 0.98 (0.90--1.07) in my depths and 0.95 (0.82--1.10) at the 817 points of \textcite{blewett2016} that I located on detected hollows (Section~S9), so the smaller depths of poleward hollows come mostly from their smaller size.

\begin{table}[!ht]
    \mytablestyle
    \caption{{\bf Temperature dependence of hollow depth compared with diffusion-limited sublimation fronts.} $\Delta H$ is the enthalpy of sublimation or vaporization per mole of vapor (individual species and sources in Table~S16). The depth ratio is the depth of a front at 324~K divided by its depth at 219~K, the range of modeled mean annual surface temperature at the measured hollows (Equation~\ref{eq:ratio}). For the measured hollows, $\Delta H$ is the apparent enthalpy implied by the fitted dependence of depth on temperature (Section~S13), for the depths of \textcite{blewett2016} and for my automatic depths.}
    \footnotesize{
        \begin{tabularx}{\textwidth}{>{\raggedright\arraybackslash}X >{\raggedright\arraybackslash}p{3.4cm} >{\raggedright\arraybackslash}p{4.4cm}}
            \toprule
            \textbf{Volatile} & \textbf{$\Delta H$ (kJ~mol$^{-1}$)} & \textbf{Depth ratio, 324~K to 219~K} \\
            \hline
            Measured hollows, depths of \textcite{blewett2016} (apparent) & 4.8 (1.6--8.1) & 1.5 (1.2--2.1) \\
            Measured hollows, this study (apparent) & $-$1.8 ($-$4.7 to 1.2) & 0.86 (0.66--1.12) \\
            Water ice & 50 & 85 \\
            Sulfur & 75--129 & 794 to 9 $\times$ 10$^{4}$ \\
            KCl and NaCl & 222--230 & 4--7 $\times$ 10$^{8}$ \\
            K, Na, Fe, Mg, and Ca sulfides & 248--477 & 4 $\times$ 10$^{9}$ to 3 $\times$ 10$^{18}$ \\
            Graphite & 717 & 5 $\times$ 10$^{27}$ \\
            \hline
        \end{tabularx}}
    \label{tab:fronts}
\end{table}

\begin{figure}[!ht]
    \centering
    \includegraphics[width=15cm]{figures/fig5_depth_v3.png}
    \caption{{\bf Hollow depth increases with hollow size and varies little with temperature and host crater size.} Open blue symbols are the shadow-length depths of \textcite{blewett2016}; filled orange symbols are my automatic depths; gray dots are my hand measurements. (a) Median depth in bands of absolute latitude against band-median mean annual temperature, with 95\% intervals from bootstrap resampling of images. The depths of \textcite{blewett2016} are corrected to their median pixel scale and solar incidence, and mine to 10~m/pixel, 70$^\circ$ incidence, and a hollow 0.5~km across. The dashed line is a sulfur sublimation front limited by vapor diffusion ($-$4,514~K), anchored at the band of \textcite{blewett2016} at 0--30$^\circ$. (b) Depth against hollow diameter for hollows $\leq$ 1~km across, corrected to 10~m/pixel and 70$^\circ$ incidence: automatic depths in diameter bins (median and 95\% interval) with the fitted power law, and the hand measurements. (c) Median depth per host crater with at least three measured hollows against host diameter, corrected to the median pixel scale, incidence, and latitude, and for my depths also to a hollow 0.5~km across; symbol area scales with the number of hollows, and the lines are the regressions on log diameter.}
    \label{fig:depth}
\end{figure}

Depth varies little with host size, nesting, and degradation state once hollow size is accounted for. Across host diameters of 5--790~km, the median depth per host crater does not change in either set (Figure~\ref{fig:depth}c; slope of log depth on log diameter $-$0.00 $\pm$ 0.04 for 87 host craters and 0.02 $\pm$ 0.06 for 56). For individual hollows, with hollow size controlled, my depths increase by a factor of 1.13 (1.02--1.25) per tenfold increase in host diameter. Hollows in host craters that formed inside older craters are as deep as hollows in isolated host craters (ratio 1.02, 0.92--1.14, in the depths of \textcite{blewett2016}; 1.01, 0.93--1.09, in mine with hollow size controlled; Section~S14). Hollows in fresh hosts are as deep as hollows in degraded hosts (1.02, 0.89--1.18; 1.01, 0.88--1.16). In the depths of \textcite{blewett2016}, hollows in hosts of Classes~3 and 4--5 of \textcite{kinczykMorphologicalEvaluationCrater2020} are 1.32 (1.14--1.51) and 1.27 (1.12--1.44) times as deep as hollows in hosts of Classes~1--2, 36 of whose 41 hosts are filled craters. My depths give similar ratios, 1.21 (1.04--1.39) and 1.30 (1.10--1.53), but the hollows in hosts of Classes~1--2 are smaller (median 340~m against 529 and 625~m), and with hollow size controlled the ratios are 1.03 (0.90--1.17) and 0.99 (0.85--1.15) (Section~S11).

Depth also varies little with composition. Across the 10th--90th percentile range of the S/Si, Mg/Si, Ca/Si, and Al/Si ratios of the MESSENGER X-ray maps \parencite{nittler2020}, depth changes by at most 11\% in the depths of \textcite{blewett2016}, and every interval includes no change; in my depths, the factor across the range of S/Si is 0.93 (0.77--1.14). Hollows where at least half of the surface within 10~km is mapped as concentrated low-reflectance material by \textcite{klima2018} are 1.09 (1.00--1.20) times as deep as other hollows in the depths of \textcite{blewett2016} and 0.99 (0.89--1.10) times as deep in mine. The depths of \textcite{blewett2016} decrease by 7\% (1--12\%) per standard deviation increase in the redness of the host crater's interior in the MDIS color map, but mine do not (factor 1.02; 0.95--1.10; Section~S9).

Depth differs between host craters in both sets (one-way analysis of variance over the 87 and 56 host craters with at least three measured hollows, $p <$ 10$^{-70}$ and $p$ = 1 $\times$ 10$^{-6}$), but craters less than 600~km apart are no more alike in depth than distant craters ($r$ = 0.01 in the depths of \textcite{blewett2016}), so depth is set locally, at the scale of a host crater or smaller (Section~S10).

\FloatBarrier
\subsection{Degradation State and Insolation Account for the Largest Shares of Variation in Hollow-Hosting Probability}
\label{sec:joint}
Among the physical factors tested in Sections~\ref{sec:source}--\ref{sec:depth}, the degradation state of a crater and its insolation account for the largest shares of the variation among craters in whether they host hollows. The tests in those sections each examine one factor, with crater size and latitude, and in most tests completeness, as controls. To compare the factors on a common scale, I fit one model of hosting that includes all of them, for the 31,554 craters of the containment tests (Section~S16).

The model approximates the hollows that form in a crater as scattered independently at random over its area $A$, with an areal density $\lambda$, the expected number of hollows per unit area, that depends on the crater's properties. The number of hollows in the crater then follows a Poisson distribution with mean $\lambda A$. Each hollow is detected with a probability proportional to $C^{g}$, where $C$ is the completeness at the crater center and $g$ is a fitted exponent, so the number of hollows detected follows a Poisson distribution with mean $\lambda A C^{g}$. The probability that a crater is observed to host hollows is the probability that at least one hollow is detected,
\begin{equation}
    P = 1 - \exp\left(-\lambda A C^{g}\right).
    \label{eq:thinned}
\end{equation}
The crater's properties enter through the logarithm of the density, $\ln\lambda = \beta_0 + \sum_k \beta_k x_k$, where each $x_k$ is a variable from Sections~\ref{sec:source}--\ref{sec:depth} and $e^{\beta_k}$ is the factor by which the density of hollows changes per unit of $x_k$. Equation~\ref{eq:thinned} is then a binomial regression with a complementary log-log link, $\ln[-\ln(1-P)] = \ln A + g \ln C + \beta_0 + \sum_k \beta_k x_k$, in which the coefficient of $\ln A$ is fixed at 1, the constant of proportionality in the detection probability is absorbed into $\beta_0$, and $g$ is fitted with the $\beta_k$. I call this the thinned occurrence model, because detection thins the hollows present in each crater, as in a thinned Poisson process, and the model predicts whether hollows occur. The fitted $g$ is 0.95, close to the value of 1 expected if the probability of detecting each hollow is proportional to completeness.

I grouped the variables by the factor they represent (Section~S16). The degradation-state group comprises the three states of \textcite{herrick2018}, their category of filled craters, and the presence of rays; I used these because the classes of \textcite{kinczykMorphologicalEvaluationCrater2020} cover only craters $\geq$ 40~km. The insolation group comprises absolute latitude, $\Delta T_\mathrm{w}$, and the change in the effect of $\Delta T_\mathrm{w}$ with crater size; these terms are correlated, but because they are in one group, the correlation does not change the group's share. The crater-size group comprises the diameter, controlled for the increase in area with diameter. The nesting group comprises whether a crater formed inside an older crater, its position on the parent's peak or floor, whether the parent hosts hollows or is a basin, and the depth of origin of the material the crater exposes. The volcanic-cover group comprises whether a crater lies on the northern smooth plains, whether it lies inside Caloris with excavation above or below the base of the plains, and whether it excavates only the volcanic fill of its parent. The vent group records whether the crater or its parent contains a volcanic vent. The composition group comprises the fraction of the crater mapped as concentrated low-reflectance material and the first two principal components of the MDIS color of its interior (Section~S9), with an indicator for the 142 craters without color. I did not include the elemental ratios of the MESSENGER X-Ray Spectrometer (XRS) in this group, because each value in their maps averages over hundreds of kilometers (Section~S9).

McFadden $R^2$ is one minus the ratio of the log-likelihood of the model to that of a baseline model \parencite{mcfadden1974}, so 0 means the model predicts no better than the baseline and 1 means it predicts perfectly. Here the baseline has completeness and crater area alone, and against it the seven groups give a McFadden $R^2$ of 0.23 (0.20--0.27). \textcite{mcfadden1977} described values of 0.2--0.4 against a baseline with a constant alone as an excellent fit. $R^2$ against that baseline is larger than against completeness and area, so 0.23 is a good fit by this measure. Dominance analysis divides this $R^2$ among the groups \parencite{azen2003,azen2009}. A group's share is the increase in $R^2$ from adding it to a model, averaged over all combinations of the other groups, as a percentage of the total $R^2$. Degradation state accounts for 37\% (29--47\%) of the explained variation and insolation for 18\% (11--21\%), followed by nesting at 16\% (12--19\%), crater size at 12\% (9--20\%), volcanic cover and composition at 7\% each (1--10\% and 3--14\%), and vents at 3\% (2--7\%), with 95\% intervals from bootstrap resampling of regions about 640~km across. Vents and the plains change the odds of hosting severalfold (Sections~\ref{sec:source} and \ref{sec:vents}), but few craters contain vents or lie on the plains, so they account for little of the variation across all craters.

On the 27,810 craters that the XRS maps cover, adding the S/Si, Mg/Si, Ca/Si, and Al/Si ratios to the composition group raises its share from 8\% to 9\% and leaves the order of the other groups unchanged (Section~S16).

The regional hosting fraction, the fraction of craters within 500~km of a crater that host hollows, measures how common hollows are in the crater's surroundings. If the seven groups captured every regional influence on hosting, adding this fraction would not improve the fit. Adding it marginally raises the McFadden $R^2$, by 0.03, so craters host hollows more often where their neighbors do, beyond what the seven groups explain. Color is already in the composition group, and adding the XRS ratios reduces this gain by only 0.001 (Section~S16), so neither the color nor the elemental composition of the surface explains the regional clustering of hollows. A hollow source that is distributed unevenly at depth and sampled only partly by the surface could produce such clustering (Section~\ref{sec:model}).

To find where the clustering lies, I compared the number of craters within 500~km of each point that host hollows with the number that the model with all seven groups predicts (Figure~\ref{fig:regional}; Section~S16). Three regions have more host craters than predicted at a family-wise significance level of 0.05, and no region has significantly fewer. In the hollow-rich province of Section~\ref{sec:sun}, 23 craters within 500~km of 65$^\circ$N, 125$^\circ$E host hollows, where the model predicts 5.0, and many of these hosts are small craters to which the model gives a low probability of hosting. I inspected all 81 hollows in these 23 craters: 79 are clearly hollows, and each of the other 2 lies in a crater with two other hollows, so all 23 craters remain hosts. The excess does not come from finer imaging: among craters 5--30~km across at 45--75$^\circ$N whose best image is at 20~m/pixel or finer, 13 of 143 in the province host hollows (9.1\%), against 9 of 831 elsewhere (1.1\%; Section~S16). Within 500~km of 13$^\circ$N, 329$^\circ$E, an area that includes the craters Abu Nuwas, Asvaghosa, Sinan, and Ts'ai Wen-chi, 23 craters host hollows, where the model predicts 8.8. Within 500~km of 7$^\circ$N, 203$^\circ$E, near Tyagaraja, 19 craters host hollows, where the model predicts 6.7; the model predicts that Tyagaraja itself hosts hollows, so the excess lies in the craters around it. Within 500~km of the points in these three regions lie 112 of the 446 host craters, against 49 predicted.

\begin{figure}[!ht]
    \centering
    \includegraphics[width=15cm]{figures/fig6_regional_excess_comp_v1.png}
    \caption{{\bf Three regions have more host craters than the joint model predicts.} Ratio of observed to predicted host craters within 500~km of each point, with predictions from the model with all seven groups (Section~\ref{sec:joint}), over a hillshade of the 665~m/pixel topography \parencite{becker2016}. Areas where fewer than two host craters are predicted are not colored. Black outlines enclose excesses at a family-wise significance level of 0.05 (Section~S16). Dots mark host craters, and the dashed line marks the rim of the Caloris basin \parencite{orgel2020}. The hot and warm longitudes are labeled at the top. (1) The hollow-rich province. (2) The area of Abu Nuwas, Asvaghosa, Sinan, and Ts'ai Wen-chi. (3) The craters around Tyagaraja.}
    \label{fig:regional}
\end{figure}
\FloatBarrier

\section{Discussion}
\label{sec:discussion}

\subsection{A Model of Hollow Formation}
\label{sec:model}
The results constrain where the hollow-forming volatile comes from, how it reaches the near surface, what drives its loss, and what sets the depth of a hollow. I propose a six-step model of hollow formation that accounts for these results (Figure~\ref{fig:lifecycle}). Table~\ref{tab:constraints} lists the observations the model must explain, the evidence for each, and the step that explains it. Below, I use ``regolith'' for fragmented rock at all depths, from the gardened surface layer to the impact-fractured rock beneath crater floors and in central peaks, and ``breccia'' only for fragmented rock bound by impact melt.

\textbf{Source step.} I propose that the ultimate source of the hollow-forming volatile is sulfide dispersed through Mercury's older crust and potentially upper mantle, and that the volatile itself is native sulfur released from that sulfide. Mercury's surface contains 2--4~wt\% sulfur \parencite{nittler2011,weider2015} and about 1.5~wt\% iron, mostly in sulfides and metal \parencite{weider2014}. Iron of 1.5~wt\% bound entirely in FeS would hold 0.86~wt\% sulfur, so most of the sulfur is in other sulfides. From X-ray spectrometer measurements, \textcite{weiderChemicalHeterogeneityMercurys2012} inferred Ca, Mg, and Fe sulfides such as oldhamite. Under the reducing conditions these compositions imply, silicate melts dissolve sulfur at weight-percent levels and crystallize Ca, Mg, and Fe sulfides \parencite{namur2016,anzures2020,boukare2019}. Low-reflectance material is enriched in graphite \parencite{peplowski2016,klima2018,vanderkaaden2015}, and linear spectral models of hollows include MgS, CaS, and graphite \parencite{emran2025}. The hollow source is closely linked to low-reflectance material, but the two need not be the same material. Composition, measured by low-reflectance material and color, accounts for 7\% of the explained variation in hosting, and host craters cluster regionally beyond what composition, including the elemental ratios, explains (Section~\ref{sec:joint}). A source distributed unevenly at depth, such as sulfide-bearing magmas from the deep crust or upper mantle, would produce this clustering, and low-reflectance material could then correlate with the source without being the source. Hollow spectra also include silicates, probably pyroxene of the host rock \parencite{lucchetti2018}. In the normative mineralogies of \textcite{mccoy2018}, calculated from MESSENGER geochemical data, two plains terranes contain the least sulfide, 3.7~wt\% each: the plains inside the Caloris basin and the Low-Fast Terrane, which is named for its low count rate of fast neutrons and lies largely within the lower latitudes of the northern smooth plains. The Northern Terrane, poleward of about 65$^\circ$N, contains 4.3~wt\% sulfide, close to the 4.7~wt\% of the southern hemisphere. Relative to the surrounding terrain, hollow fields on the plains are nonetheless scarcer poleward of 60$^\circ$N, 0.13 times as dense as outside the plains at 60--80$^\circ$N against 0.35 times at 40--60$^\circ$N (Section~\ref{sec:source}). The sulfide content of the plains therefore does not explain the scarcity of hollows on them. Sulfides themselves, however, are too refractory to sublimate at Mercury's surface temperatures \parencite{phillips2021}, and native sulfur forms from sulfide only where another element is reduced. In the model, impact heat first converts sulfide to native sulfur, reducing iron to metal or, in impact melt, silica to SiO and SiS gas (charging step). This crust lies beneath young volcanic plains, so hollows form where impacts excavate it: hollows are scarce on the plains, and 19 of the 20 craters on the plains that host hollows are $\geq$ 10~km across (Section~S6); hollow fields in Caloris occur only in craters whose excavation reaches below the plains; and hollows are densest on central peaks and peak rings (Section~\ref{sec:source}). Mercury's crust is 35 $\pm$ 18~km thick on average \parencite{padovan2015}, so the largest craters and basins expose rock from below the crust: 69 of the 532 host craters in the crater containment graph have a first-order depth of origin deeper than 35~km (Section~S14). Sulfide-rich layers may also form in Mercury's mantle as its magma ocean solidifies \parencite{boukare2019}, so part of the sulfide source may have originally come from the mantle. The chloride spectra of hollows in Tyagaraja crater \parencite{lucchetti2021} and the stability of dark spots \parencite{wang2022} remain to be reconciled with a sulfide source. The model here focuses on sulfides that release sulfur and sulfur-bearing gases, but hollow-forming impacts on Mercury may release a mixture of gases, which could contribute to the carbon- and chloride-bearing phases identified in hollows \parencite{emran2025,lucchetti2021}.

\textbf{Charging and collapse step.} Impacts expose the older crust and heat it. The rock that rises into a central peak comes from just below the maximum depth of melting \parencite{cintala1998}, so it is the most strongly heated rock that the impact leaves unmelted. Hollows are densest on peaks (Section~\ref{sec:source}), so in the model the charged rock lies mainly over the most intensely heated regions of the impact. In equilibrium with iron metal, the S$_2$ pressure over FeS is 0.1~Pa at 1,344~K and 1~Pa at 1,471~K \parencite[Section~S13;][]{janaf_fe025}. Where its vapor can escape, FeS loses sulfur and leaves porous metallic iron \parencite{tachibana1998}. The breakdown is an internal redox reaction: the oxidation of sulfide to native sulfur is balanced by the reduction of Fe$^{2+}$ to metal, so no external oxidant is needed \parencite{brett1976}. Ca and Mg sulfides release S$_2$ where silica is reduced in place of iron. Silicon vaporizes from silicate melts mainly as SiO gas \parencite{schaeferCHEMISTRYSILICATEATMOSPHERES2009,sossiThermodynamicsElementVolatility2018}, and the reduction of silica has been proposed on Mercury before, by graphite in lavas \parencite{mccubbinLowOSi2017}. Experiments run the same chemical system in the opposite direction: reducing, sulfur-rich gas converts glass of Mercury's surface composition to Ca and Mg sulfides and silica \parencite{renggli2022}. For MgS + SiO$_2$ $\rightarrow$ MgO + SiO(g) + $\frac{1}{2}$S$_2$(g), the equilibrium S$_2$ pressure over pure phases is 1~Pa at 1,600~K, 26~Pa at 1,800~K, and 7~kPa at 2,300~K; for the corresponding reaction of CaS, it is 0.02, 0.7, and 460~Pa (Section~S13; thermodynamic data of \cite{mcbride2002}). With the same data, the S$_2$ pressure over FeS is 44~Pa at 1,800~K and 2.3~kPa at 2,300~K. Silicon can also vaporize as SiS: for 2MgS + SiO$_2$ $\rightarrow$ 2MgO + SiS(g) + $\frac{1}{2}$S$_2$(g), the S$_2$ pressure is 87~Pa at 1,800~K, and forming enstatite instead of MgO raises it further, so the pressures for MgS are lower limits; for CaS, the reaction that forms SiO releases more sulfur. In the model, Ca and Mg sulfides therefore release sulfur only where rock exceeds about 1,800~K, in impact melt and in the most strongly heated rock beside it. As the gas cools, SiO condenses as silica and silicon at 1,300--1,500~K for SiO pressures of 1--50~Pa, whereas S$_2$ at the same pressures condenses at 500--600~K (Section~S13). Iron remains as metal in the hot rock, so the sulfur that reaches cooler rock condenses as native sulfur nearer to the surface.

Impact melt supplies sulfur to charged rock only where the melt is enclosed. Gas that leaves the surface of a melt pond on a crater floor escapes to space, because no rock lies above the pond in which the vapor could condense. Melt injected as dikes and veins into the floor and the central uplift, and melt-bearing breccia buried beneath fallback or slumped wall material, release their gas into the surrounding rock. Within bins of distance from the crater center, floor within 0.1 radii of the walls or the peak has 1.04 (0.85--1.27) times the hollow density of floor farther from them, floor beside the walls alone has 0.87 (0.69--1.08) times, and floor beside the peak alone has 1.18 (1.00--1.39) times (Section~S4). Peak terrain within 2~km of the floor has half the hollow density of the rest of the peak (Section~\ref{sec:source}). Hollows are therefore no denser where a melt sheet meets the walls and only slightly denser where it meets the peak, so sulfur from an open melt sheet is not needed to explain where hollows occur.

The vapor moves along impact fractures. Where the impact-melt breccia is still hot, the vapor converts glass to Ca and Mg sulfides with silica and aluminosilicate, which leaves the rock porous and fragile \parencite{renggli2022,renggli2023}. Vapor at 0.1--1~Pa condenses as native sulfur only where the regolith is cooler than about 340--375~K (Section~S13): in the fractured walls, rim, and slumped material of the crater, and in regolith above and beside the impact-heated rock. I call regolith that contains this native sulfur the charged layer; sulfidized regolith lies beneath and beside it, near the impact-heated rock and along the hottest fractures. Where the vapor spreads laterally through fractured regolith, the charged body has the irregular outline of an individual hollow; hollow fields form as neighboring hollows widen by scarp retreat and merge.

Native sulfur that fills pores and fracture throats as it condenses seals them, and heat from impact-heated rock then warms it in place. In a sealed pore, the pressure of sulfur vapor equals the vapor pressure of the sulfur. Sulfur melts at 393~K, the triple point of monoclinic sulfur, liquid, and vapor, where its vapor pressure is 2.4~Pa \parencite{steudel2003}. Above this temperature, sulfur in a sealed pore is liquid in equilibrium with its vapor. Exposed sulfur, solid or liquid, is lost by evaporation. Sulfur remains in a sealed pore only while its vapor pressure is below the confining pressure, which is at most the weight of the overlying column, $\rho g z$ at depth $z$, where $\rho$ is the bulk density of the column and $g$ is gravity. Where the regolith becomes hotter than the boiling temperature of sulfur under that load, the pores fail, and the vapor rises and recondenses in cooler regolith above. Where the peak temperature rises with depth faster than the boiling temperature does, 0.4--1.6~K per meter at the depths in Section~\ref{sec:depthcontrol}, the charged layer therefore has a base at the depth $d$ where the two are equal, $P_{\mathrm{S}}(T) = \rho g d$, with $P_{\mathrm{S}}$ the vapor pressure of sulfur and $T$ the peak temperature at $d$; native sulfur remains above $d$ and none below. Where the peak temperature is nearly uniform with depth, as in regolith heated from the side, the order reverses: sulfur boils above $d$ and remains below it. With tensile strength, the base would lie deeper, and with confinement by a horizontal stress smaller than the overburden, the base temperatures would be 40--120~K lower (Section~S13).

The depth relation requires native sulfur to be in place before the regolith reaches its peak temperature. Single S$_2$ molecules cross 100~m of fractures 10~$\mu$m to 1~mm wide in 0.6--60 days, but at the S$_2$ pressures of FeS the vapor transports little sulfur, and the front of regolith filled with 5~vol\% sulfur advances 8--110 times more slowly than the heat front (Section~S13). Sulfur that the same impact releases from FeS therefore condenses mostly beyond the heated regolith or in regolith that has already passed its peak temperature, where its vapor pressure stays low. This sulfur is stored as native sulfur, solid or liquid, in pores; it moves again as vapor only where later heating raises its vapor pressure above the confining pressure, and the pores it leaves can then compact. Native sulfur is in place before the heat arrives where an impact heats the charged layer of an older crater, as the nesting results imply for many host craters (Section~\ref{sec:nesting}). Hollows in nested and isolated host craters have the same depth, however (Section~\ref{sec:depth}), so the process that sets depth acts in both. For isolated hosts, the model therefore requires sulfur delivered faster than diffusion from heated FeS can supply it. In impact melt enclosed in the floor or the uplift, MgS releases sulfur at S$_2$ pressures of at least 1~Pa at 1,600~K and 7~kPa at 2,300~K, and gas at these pressures moves through fractures by pressure-driven flow, the process that could place native sulfur in cooler regolith before the heat arrives. Shock heating of sulfide during excavation, or native sulfur already present in the target crust from earlier, smaller impacts, could also supply it. Isolated craters with explosive vents host hollows 2.51 times as often as other isolated craters (Section~\ref{sec:vents}), which fits intrusions as a second source of heat or sulfur (Figure~\ref{fig:lifecycle}b, lower pane). Explosive eruptions require magma rich in volatiles \parencite{kerber2011}, so the association also fits a source region rich in volatiles. Hollows in craters with vents are as deep as hollows elsewhere, which suggests that in these craters, too, the depth was set where the vapor pressure of sulfur balanced the overburden at a base temperature reached by impact heating.

Collapse of the charged layer begins while the crater is hot, and in this model it produces the depth of a hollow. Hot sulfur vapor corrodes the regolith it moves through, and its reactions with natural regolith are likely more varied than the sulfidation measured in experiments \parencite{renggli2022,renggli2023}. At the base of the charged layer, the pressure of sulfur vapor in the pores equals the confining pressure, so the effective stress on the grain framework is near zero. Where pores fail and their vapor escapes, the full weight of the column rests on the weakened grain framework, which compacts. Native sulfur at the 820--1,050~K of the base (Section~\ref{sec:depthcontrol}) has 7 $\times$ 10$^6$ to 3 $\times$ 10$^{10}$ times its vapor pressure at present mean annual temperatures of 270--330~K (Section~S13), so sulfur in pores that reopen while the crater is hot, for example as the regolith cools and contracts, leaves quickly. As vapor escapes, the lower part of the charged layer loses mass and pore pressure, and the column above it compacts. If the charged layer loses its sulfur and compacts while the crater is hot, the depth of a hollow forms during charging and does not depend on present-day temperature. Present-day activity, driven by sunlight and fed by native sulfur in the cooler upper part of the layer, is confined mainly to scarps and floors and keeps hollows fresh. Neither the reopening of the pores nor the compaction has been modeled.

On Earth, ground subsides in the same way where mass is lost from below. In karst terrain, groundwater dissolves rock and removes it in solution, and the ground above collapses into sinkholes \parencite{ford2007}. In thermokarst terrain, ground ice melts and the ground subsides into depressions with steep walls and flat floors, which widen as their walls retreat and merge with one another \parencite{czudek1970}; hollows have the same form. On Mercury, the lost mass leaves only as vapor.

\textbf{Storage step.} Native sulfur sublimates from the top decimeters and leaves a sulfur-free lag above the charged regolith \parencite{blewett2013,blewett2016,thomas2016}. Beneath a lag 0.13--15~m thick, sulfur loss slows to 3.3--6.4~m per Gyr \parencite{phillips2021}, so a charged layer persists long after the impact heat has dissipated. Because native sulfur persists beneath a lag, hollows can be active in host craters whose impact heat dissipated more than 10$^8$~yr earlier (Section~\ref{sec:age}). Sulfur that returns from the exosphere also diffuses into porous regolith and collects in modeled reservoirs 8--21~cm deep, whose longitudinal distribution resembles that of hollows \parencite{verkercke2025}. Exospheric implantation could therefore add to, and replenish, the top of a charged layer, contributing to present-day activity.

\textbf{Activity step.} Small impacts, mass wasting, and scarp retreat expose the charged layer \parencite{blewett2016,wright2020,speyerer2022,deutsch2025}; in the most degraded host craters, hollows lie mostly in small superposed craters \parencite{giroud2024}. Sulfur supplied from below condenses in the exposed regolith at night, and crystals that grow in pores may push grains apart, as salt crystals do in salt weathering \parencite{scherer2004,melosh2011}. During the day, the top 10--22~cm of an equatorial site heats past 368~K (inset of Figure~\ref{fig:lifecycle}e). At 368~K, native sulfur changes from the orthorhombic to the monoclinic form and expands by about 3\%, from the densities of the two forms at 298~K \parencite{steudel2003}, and at higher temperatures it evaporates. The expansion pushes grains apart, and grains settle into the pore space left as the sulfur evaporates, so the regolith compacts as it loses its sulfur. On a scarp, the loosened regolith falls away and exposes fresh charged regolith, so the scarp retreats \parencite{blewett2016}. The sulfur-free debris, $1-\phi$ of the volume, remains, so a hollow widens at constant depth only where the charged column beneath the retreating scarp collapses as it loses its sulfur or where the debris is removed. On slopes, debris can move away: in MESSENGER repeat images, bright material moved down a crater wall from a hollow in Sholem Aleichem crater within six months in 2011 \parencite{speyerer2022}. On a flat floor, the compacted, sulfur-free residue is $(1-\phi)/\phi$ times as thick as the lowering of the floor, where $\phi$ is the fraction of the layer's volume that native sulfur and its pore space occupied. After 1--15~cm of lowering for $\phi$ = 0.1--0.4, the residue is thicker than the layer that crosses 368~K each day, so the daily breakup no longer reaches charged regolith, although sulfur still diffuses slowly through the residue (Section~S13). In Dominici crater, vertical growth appears to have stopped at a silicate-dominated residue \parencite{emran2025}. A charged layer tens of meters thick beneath a flat floor therefore cannot collapse by daily cycling alone, so the full depth of a hollow forms only by collapse while the crater is hot (charging and collapse step) or where the residue is removed. If the crater-retention age of hollow floors, $\leq$ $\sim$10$^5$~yr, dated whole hollows, hollows would deepen at $>$ 0.23~mm per yr \parencite{wang2020}; if the depth formed during charging, that age dates the renewal of the floor surface. Native sulfur cannot persist on sunlit ground, where 1~mm of it evaporates in about 2 days at 300~K (Section~S13), so the bright halo more likely consists of products of the reaction between sulfur vapor and the surrounding regolith, as sulfur-rich gas reacts with Mercury surface materials in experiments \parencite{renggli2022}, and the MgS and CaS in hollow spectra \parencite{vilas2016,emran2025} may include such products.

Sunlight modulates this step. Craters near the hot longitudes host hollows more often near the equator and to 30$^\circ$N, the odds of hosting rise with $\Delta T_\mathrm{w}$, and at high northern latitudes hollows are more common on equator-facing walls (Section~\ref{sec:sun}); \textcite{deutsch2025} likewise concluded that insolation drives the initiation and growth of hollows. Fresh, Stage 1 hollows make up nearly all hollows where $\Delta T_\mathrm{w}$ exceeds about 50~K, so the most vigorous activity may require a threshold of insolation, perhaps implying a typical lag thickness (Section~\ref{sec:sun}). Sunlight does not limit where hollows occur, however: hollows are common at the warm longitudes and at high northern latitudes, hollow depth changes little with temperature (Section~\ref{sec:depth}), and in the joint model insolation accounts for 18\% of the variation in crater hosting that physical factors explain, after degradation state (Section~\ref{sec:joint}).

\textbf{Dormancy step.} Activity stops where the exposed charged regolith is spent, where scarps retreat into regolith without native sulfur, or where a lag covers the scarps and floor. The halo then fades; fading turns a bright, sharp Stage 1 hollow of \textcite{deutsch2025} into a Stage 2 hollow of the same size and shape. Exposed charged layers are spent or buried over time, so active hollows should be more common in young host craters, as observed (Section~\ref{sec:age}). If collapse during charging sets the depth, hollows of a given size should also be as deep in young hosts as in old ones. Hollows in fresh and degraded hosts of \textcite{herrick2018} have the same depth, and hollows in hosts of Classes~1--2 of \textcite{kinczykMorphologicalEvaluationCrater2020} are shallower than hollows in hosts of Classes~3--5 because they are smaller; with hollow size controlled, their depths do not differ (Section~\ref{sec:depth}). Gardening also fills old hollows, but only slowly (fossilization step). Small impacts that strip the lag can return a hollow to activity \parencite{blewett2016}.

\textbf{Fossilization step.} Gardening and space weathering soften the scarps and erase the reflectance contrast, which produces Stage 3 hollows. Stage 3 hollows are rare, 10 of 2,666 classified hollows \parencite{deutsch2025}, although topographic diffusion lowers relief by only 0.4--2.7~m per Gyr \parencite{phillips2021}. Their rarity implies faster infilling, frequent return to activity, or a bias against detecting hollows without halos (Section~S12).

\begin{figure}[!htbp]
    \centering
    \includegraphics[width=13cm]{figures/fig7_lifecycle_v12.png}
    \caption{{\bf The life cycle of a hollow.} The schematic is not to scale; timescales are given in parentheses, and Section~S12 lists their bases. Panels a and c--g show the six steps of hollow formation, numbered 1--6, and panel b shows the heat trigger for context. (a) Step 1, source (buried $\geq$ 2.5~Gyr ago): older crust or upper mantle with dispersed sulfides, linked to low-reflectance material, lies beneath younger volcanic plains. (b) Heat trigger (at the scale of the crater): an impact excavates through the younger plains into the older crust and spreads ejecta of older crust over the plains; the central uplift raises rock from just below the maximum depth of melting into the central peak, impact melt lines the floor, and impact-heated rock lies beneath it. The box marks the part of the crater floor that c--g enlarge, and the hot zone in c is the impact-heated rock beneath the floor. The pane below it shows a possible second source of heat or sulfur for isolated host craters: magma intruded beneath the crater floor, drawn as a sill that also feeds an explosive vent. (c) Step 2, charging and collapse (minutes to 10$^6$~yr): an impact exposes and heats the older crust; above $\sim$1,300~K, FeS releases S$_2$ and leaves iron metal, and above $\sim$1,800~K, in impact melt, Ca and Mg sulfides release S$_2$ as silica is reduced to SiO and SiS gas. The S$_2$ rises along fractures, sulfidizes the regolith near the impact-heated rock, and condenses as native sulfur in cooler regolith above the depth $d$ where the vapor pressure of sulfur at the regolith's peak temperature equals the overburden. Where hot vapor reacts with the lower part of the layer and escapes from it, the porous, weakened regolith compacts under the weight of the overlying column, and the surface sags. (d) Step 3, storage (up to $>$ 10$^8$~yr): native sulfur remains in the cool upper part of the layer beneath a sulfur-free lag; compacted residue fills the lower part. (e) Step 4, activity (floors $\leq$ $\sim$10$^5$~yr old; scarps retreat $\geq$ 1~cm per 10$^4$~yr): where the charged layer is exposed, sulfur condensed at night expands at 368~K through the orthorhombic-to-monoclinic phase transition and sublimates by day, and the loosened regolith falls from the scarp, which retreats into charged regolith at the margin of the charged body. The floor lies at depth $h = \phi d$ above compacted residue; products of the reaction between sulfur vapor and the surrounding regolith form a bright halo. (f) Step 5, dormancy ($\sim$0.7 times as long as activity; halos fade within $\sim$280~Myr): a lag covers the scarp and floor, and the halo fades. (g) Step 6, fossilization ($>$ 4~Gyr to erase the relief by diffusion alone): gardening and space weathering soften the hollow's appearance. Inset of e, at the top of the scarp (dashed box in e): daily maximum (solid) and minimum (dashed) temperature with depth at the hot (orange) and warm (blue) equator from the thermal model of Section~\ref{sec:craters}. Horizontal lines mark the base of the layer that crosses 368~K every Mercury solar day, and the shading above them marks the zone of activity, where present-day activity is highest in this model.}
    \label{fig:lifecycle}
\end{figure}
\FloatBarrier

\subsubsection{Hollows in Nested Craters and on the Caloris Knobs}

\textbf{Craters inside charged craters.} Craters inside older craters host hollows more often, and craters inside hollow-bearing parents host hollows more often still (Section~\ref{sec:nesting}). In the charging and collapse step, native sulfur condenses in the walls and rim of a crater as well as beneath its floor and in its central peak, so a crater on the parent's wall exposes charged regolith nearly as often as a crater on its peak. The excavation depth of a crater $\geq$ 5~km is at least 0.4~km, so the impact that forms a crater inside a charged parent cuts through the parent's charged layer, 60--280~m thick for $\phi$ = 0.1--0.4 (Section~\ref{sec:depthcontrol}), and the walls and rim of the new crater expose that layer. Its impact also heats the parent's native sulfur in place, which meets the timing requirement of the charging and collapse step. Craters inside parents without hollows host hollows at 0.53\%, a seventh of the 3.8\% inside hollow-bearing parents (Section~\ref{sec:nesting}). Some of these parents may once have had hollows that have since faded to Stage 3 or been erased by later impacts, so the 44 parents without observed hollows whose children host hollows may be good places to look for faded hollows, although most are poorly imaged, with a median completeness of 0.14 at their centers (Section~S14).

\textbf{Uplifted rock outside Caloris.} In basins as large as Caloris, the central uplift rises beyond a stable height and collapses. In models of Orientale, the collapse of the central uplift formed the Inner Rook ring \parencite{johnsonFormationOrientaleLunar2016}, and in models of Caloris, crust 350--700~km from the center flows outward, folds, and overturns \parencite{gosselinCrustalBlockMuted2023}. The knobs around Caloris are usually interpreted as ejecta blocks \parencite{fassett2009,ackissKnobHeightsCircumCaloris2015}. Collapse of the central uplift would also throw blocks beyond the rim, so the two interpretations differ mainly in the depth from which the blocks came. If some knobs came from below the maximum depth of melting, they are analogous to the central complexes of smaller complex craters, and the model predicts that knobs would host hollows as central peaks do. Hollows on the knobs form in the knob material \parencite{wright2020}, and the knobs have 4.0 times the hollow density of the terrain around them and about the same hollow density as the peaks of the most degraded craters (Section~\ref{sec:source}). Mass wasting on steep knob slopes also exposes fresh material \parencite{wright2020} and could raise the density of hollows on the knobs.

\subsection{Sulfur Vapor Pressure and Collapse Set Hollow Depth}
\label{sec:depthcontrol}
In the charging and collapse step, the base of the charged layer lies at $d = P_{\mathrm{S}}(T)/(\rho g)$, where $T$ is the peak temperature at the base and $g$ = 3.70~m~s$^{-2}$ on Mercury. Sulfur boils at 718~K under 0.101~MPa \parencite{lide2007}, the weight of 9--14~m of rock with bulk density 2,000--3,000~kg~m$^{-3}$ on Mercury. If the whole layer above the base loses its sulfur and pore space, the floor lowers by $h = \phi d = \phi P_{\mathrm{S}}(T)/(\rho g)$, the largest possible lowering; if native sulfur remains in a fraction $f$ of the layer, the floor lowers by $\phi(1-f)d$. For complete collapse, the median depths of 24--28~m equatorward of 60$^\circ$ in the measurements of \textcite{blewett2016}, corrected to their median pixel scale and incidence, require a charged layer 60--70~m thick and a base at 820--890~K if $\phi$ = 0.4, 120--140~m and 880--960~K if $\phi$ = 0.2, and 240--280~m and 950--1,050~K if $\phi$ = 0.1 (Section~S13). Depths measured by different observers differ by about 20\% (Section~\ref{sec:depth}), which would change the thickness of the layer by 20\% and its base temperature by 15--30~K. Figure~\ref{fig:tbasephi} gives the depth that complete collapse produces for any base temperature and $\phi$, for future comparison with thermal models of cooling craters. Among the sulfur, chloride, and sulfide candidates, only native sulfur reaches the vapor pressure that balances the overburden at these depths at temperatures below those of impact-heated rock, 1,300--1,700~K; KCl and NaCl would require 1,920--2,430~K, above their boiling points, and MgS, CaS, and FeS more than 2,700~K (Section~S13). Above the normal boiling point, these temperatures come from a Clausius--Clapeyron relation anchored at 718~K, with slopes that span a fit to the vapor pressures at 591--717~K and the chord to the tabulated critical point of sulfur \parencite[1,314~K, 20.7~MPa;][]{lide2007}.

\begin{figure}[!ht]
    \centering
    \includegraphics[width=12cm]{figures/figS_tbase_phi_v1.png}
    \caption{{\bf Hollow depth produced by complete collapse.} Hollow depth for complete collapse, $h = \phi P_{\mathrm{S}}(T_{\mathrm{base}})/(\rho g)$, against the peak temperature at the base of the charged layer and the volume fraction $\phi$ lost on collapse. Colors and contours give $h$. The band marks the median depths of 24--28~m equatorward of 60$^\circ$ in the measurements of \textcite{blewett2016}, corrected to their median pixel scale and incidence; the difference of about 20\% between observers (Section~\ref{sec:depth}) extends this range to about 20--34~m. The top axis gives the thickness of the charged layer, $d = P_{\mathrm{S}}(T_{\mathrm{base}})/(\rho g)$, which depends only on $T_{\mathrm{base}}$. $P_{\mathrm{S}}$ uses the mean of the two slopes of the vapor-pressure relation above the normal boiling point (Section~S13); $\rho$ = 2,500~kg~m$^{-3}$ and $g$ = 3.70~m~s$^{-2}$. A density of 2,000 or 3,000~kg~m$^{-3}$ moves the band about 20~K cooler or hotter, and the two slopes move it by up to 20~K.}
    \label{fig:tbasephi}
\end{figure}

In the model here, the base of the charged layer and the volume it loses set the depth of a hollow. Depth therefore depends on neither crater size nor present-day temperature if the bases of the charged layers formed by different craters reach a similar peak temperature (that is, if the peak temperature 60--280~m below the surface is similar regardless of crater size) and lose similar volume fractions on collapse, because gravity and the vapor-pressure curve of sulfur do not depend on either. The median depth of hollows does not change with host diameter from 5 to 790~km (Section~\ref{sec:depth}), whereas the sulfur supply per unit area of the floor should grow with the thickness of rock that an impact heats, which increases with crater size \parencite{melosh1989}; $\phi$ must therefore be nearly independent of the sulfur supply, so $\phi$ probably consists mostly of pore space that closes on collapse. Depth increases with hollow size, however, and fields of large hollows are deeper than their size alone predicts (Section~\ref{sec:depth}). A $\phi$ that differs between hollow fields would make the hollows of some fields deeper than others, and if regolith with more pore space and sulfur also collapses over wider areas, it would make their hollows larger as well. Within a field, depth also increases with size, which a $\phi$ that varies below the scale of a field, or hollows that deepen as they widen by scarp retreat, could produce. The shadow depths cannot separate these explanations.

The depth control rests on the peak temperature at the base of the charged layer and on $\phi$. I inferred the temperature at the base of the charged layer from the hollow depths for assumed values of $\phi$ (Figure~\ref{fig:tbasephi}); native sulfur there condensed in regolith cooler than about 375~K and was heated to 820--1,050~K after it arrived (charging step). Because the median depth does not change with crater size, this temperature must vary little among craters of all sizes: at fixed $\phi$, a twofold range of depth corresponds to a range of 60--90~K in base temperature (Section~S13). Regolith heated by contact with impact-heated rock reaches at most the contact temperature, the average of the two initial temperatures weighted by their thermal inertias, $\sqrt{k \rho c}$, which does not depend on the sizes of the two bodies while both are thicker than the distance heat has diffused \parencite{carslaw1959}. For rock at 1,300--1,700~K and regolith at 350~K, the contact temperature is 825--1,025~K for equal thermal inertias and 980--1,250~K for rock with twice the thermal inertia of the regolith, close to the range the depths imply. Heat crosses 60--280~m in 110--2,500~yr, however, after which radiation from the surface and the size of the hot zone also shape the temperature profile, so a thermal and vapor-transport model of cooling craters is needed to test whether base temperatures are as uniform as the depths require.

Two steady states could make the base temperature nearly independent of crater size: conduction above the impact-heated rock, whose contact temperature does not depend on crater size, and a two-phase zone of sulfur vapor and liquid held near the boiling curve of sulfur (Section~S13). In the first, a twofold range of hollow depth allows about a 1.8-fold range in the depth of the hot rock, and the steady state may not apply beneath craters 5--10~km across (Section~S13).

\subsection{Other Models}
\label{sec:others}
Elements of the model proposed here come from earlier models, and the model differs from each in at least one step. Exhumation of volatile-rich material followed by a protective lag \parencite{blewett2013,blewett2016} overlaps with the source, storage, and dormancy steps; the model adds charging, a depth set by collapse, and a mechanism that breaks the regolith apart. A sublimation front driven by sunlight \parencite{phillips2021} predicts a 794-fold change in depth across the measured hollows, where the observed change is 1.5-fold in the depths of \textcite{blewett2016} and 0.86-fold in mine. Exospheric sulfur \parencite{verkercke2025} is part of the model as a possible addition to the top of the charged layer. With a buried volatile-rich layer \parencite{rodriguez2020,rodriguez2023}, one might expect depth to be coherent across a region, but craters less than 600~km apart are no more alike in depth than distant craters. This lack of regional coherence does not exclude a global layer at equal depth, although such a layer would not explain why depth differs between host craters (Section~\ref{sec:depth}). A layer of sulfide that separates from impact melt and floats on it \parencite{vaughan2012}, as Ca and Mg sulfides collect in the slag that floats on molten metal in steelmaking \parencite{helbert2013}, would place hollows on melt-covered floors. Such a layer would not form on the peaks, walls, and ejecta where many hollows occur \parencite{blewett2018}. Its thickness would vary with the volume of melt and its sulfur content, whereas hollow depth does not vary with host size (Section~\ref{sec:depth}), and its sulfur would remain in refractory sulfide. With a volatile source at a fixed depth, hosting would also increase with the depth a crater reaches, but depth of origin adds little to nesting outside Caloris, and within the same parent, craters on its wall host hollows nearly as often as craters on its peak or floor (Section~\ref{sec:nesting}). A chloride-rich layer from differentiated impact melt \parencite{lucchetti2021} is consistent with the spectra of the Tyagaraja hollows but not with the scarcity of hollows on the northern smooth plains (Section~\ref{sec:source}), where the Northern Terrane contains 0.45~wt\% Cl, about three times the 0.14~wt\% of the southern hemisphere \parencite{mccoy2018}. In the collapse model, KCl and NaCl would also need base temperatures above their boiling points and above those of impact-heated rock to produce the measured depths (Section~\ref{sec:depthcontrol}). Reaction of the surface with sulfur-rich gas \parencite{renggli2022} requires a persistent supply of that gas, which impacts provide in the model. \textcite{bickel2025} attributed hollow formation mainly to micrometeoroid bombardment and insolation; in the model, micrometeoroid gardening could be one of the processes that breach the lag.

In a model in which the depth forms after exposure, as broken regolith loses its sulfur and the sulfur-free residue is removed, depth would not need to form while the crater is hot. Such a model must remove tens of meters of solids from closed, rimless depressions. Downslope transport cannot carry debris out of a closed depression, topographic diffusion lowers relief by only 0.4--2.7~m per Gyr \parencite{phillips2021}, and the gas released by volatile loss is not vigorous enough to loft dust grains \parencite{blewett2016}. Depth that increases with hollow size (Section~\ref{sec:depth}) would fit hollows that deepen as they widen, but the removal of the residue remains unexplained, so I disfavor this model.

\subsection{Untested Steps and Predictions}
\label{sec:predictions}
Five parts of the model have not been tested. First, native sulfur must be in place before the regolith reaches its peak temperature, which diffusion from FeS heated by the same impact cannot achieve (charging and collapse step); this holds for nested host craters, and for isolated hosts it requires a faster release of sulfur, such as through pressure-driven flow, or sulfur already in the target. Second, the charged layer must collapse while the crater is hot, because beneath a flat floor at present-day temperatures collapse stops once the residue forms, and the peak temperature at the base must be similar in craters of all sizes. Both conditions can be tested with a coupled model of heat conduction, vapor transport and reaction, pore failure, and compaction in cooling craters. Third, daily cycling must break sulfur-charged regolith apart on exposed scarps, and the debris must be removed or the column beneath the scarp must collapse. Thermal cycling between 82 and 293~K, which does not cross the orthorhombic-to-monoclinic transition of sulfur at 368~K, reduced the strength of sulfur-cemented lunar regolith simulant by 80\% in 80 cycles \parencite{grugel2008}, but no laboratory study has cycled native sulfur in rock across 368~K in vacuum; in the closest experiments, sulfides were heated under Mercury surface conditions \parencite{helbert2013,varatharajan2019,barbaro2023}. Fourth, enough sulfur must be released from the impact-heated rock. Heating rock with 1.5~wt\% iron, all as FeS, releases at most 0.86~wt\% sulfur, so a layer that contains 10~vol\% native sulfur requires the sulfur from at least eight times its own volume of rock heated above $\sim$1,300~K, and diffusion from FeS may deliver it fast enough only beneath large craters (Section~S13). FeS becomes unstable in very reduced melts \parencite{anzures2020}, however, and surface iron on Mercury may be mostly sulfide \parencite{weider2014,renggli2026}. Fifth, hollow floors must be renewed, by small impacts that breach the residue and by scarp retreat, fast enough to stay free of craters for $\leq$ $\sim$10$^5$~yr after the depth has formed (Section~S12). The absence of craters on hollow floors may partly reflect survivorship, however: an impact that leaves a resolvable crater may erase a hollow or hide it, and smaller craters are unresolved or are consumed as the hollow grows.

The model makes predictions that can be tested with new data. (1) Within a hollow field, depth increases with hollow size up to a common maximum, the depth of complete collapse, and this maximum depth differs between adjacent fields. (2) Hollows of a given size are as deep in the smallest host craters as in large host craters; hollows on the rims of small craters, which may have formed where the overturned flap exposed the charged layer directly, may differ. (3) Scarps retreat fastest at the hot longitudes and on equator-facing walls, and floors deepen far more slowly than scarps retreat. (4) Craters whose excavation does not reach through a plains unit of known thickness lack hollow fields on their walls, peaks, and rims. (5) Hollows and their halos expose Ca and Mg sulfides with silica and aluminosilicate, which have diagnostic emissivity features at 7--14~$\mu$m \parencite{renggli2023}. (6) Regolith that contains native sulfur loses cohesion when cycled across 368~K in vacuum. (7) Halos consist of reaction products and have no measurable relief. (8) Hollows on the knobs around Caloris are as deep as hollows of the same size in craters. (9) The three regions where more craters host hollows than the joint model predicts (Section~\ref{sec:joint}) differ in composition from their surroundings in ways that MDIS color and the X-ray maps do not resolve. Predictions 1, 2, 3, and 8 can be tested by comparing BepiColombo SIMBIO-SYS stereo and color images \parencite{cremonese2020,rothery2020} with MESSENGER images \parencite{speyerer2022,deutsch2025}, predictions 5 and 9 with the MERTIS thermal infrared spectrometer \parencite{hiesinger2020}, prediction 6 by laboratory experiments, and prediction 7 with shadows at the edges of halos in MDIS and SIMBIO-SYS images at high solar incidence.

The size distribution of hollows gives two predictions that BepiColombo can test. The High Resolution Imaging Channel of SIMBIO-SYS will image about 20\% of Mercury's surface at up to 6~m/pixel \parencite{cremonese2020}, at which a hollow 0.1~km across spans about 17 pixels. If the power law between 0.3 and 3~km continues to smaller sizes, these images will show 4.7--5.8 times as many hollows larger than 0.1~km as there are hollows larger than 0.3~km at the same sites, and about 117,000--186,000 hollows larger than 0.1~km on the surface the MESSENGER images searched here cover. Fewer would show that the distribution flattens below 0.3~km: a sharp drop at one size would point to a minimum hollow size, whereas a gradual decline would point to a process that makes small hollows rarer, such as the merging of neighboring hollows. At the large end of hollow sizes, new images of terrain MESSENGER imaged should add few hollows larger than 5~km: about 105 in all, against 96 in the catalog.

\FloatBarrier
\section{Conclusions}
\label{sec:conclusions}
I propose a six-step model for hollow formation on Mercury. Sulfides in older crust, and potentially upper mantle, are the ultimate source, and native sulfur is the volatile. Impact heat releases sulfur vapor from FeS above about 1,300~K and, in impact melt above about 1,800~K, from Ca and Mg sulfides as silica is reduced. The vapor condenses as native sulfur in cooler regolith, down to the depth where the vapor pressure of sulfur at the regolith's peak temperature balances the overburden. While the crater is hot, the vapor escapes from the lower part of this charged layer, and the porous, weakened regolith compacts under its own weight, as ground subsides where groundwater dissolves rock (karst) or permafrost thaws (thermokarst). The lowering of the surface during this compaction is the depth of a hollow. Native sulfur remains beneath a sulfur-free lag in the cooler upper part of the layer, and once the layer is exposed, nightly condensation, daily expansion at 368~K, and sublimation drive scarp retreat, which keeps hollows fresh. Global data are consistent with each step. Hollows form where impacts expose the crust beneath the volcanic plains, are densest on crater peaks, and are more common at the hot longitudes near the equator and to 30$^\circ$N, in fresh craters, in craters inside older craters, and in craters with explosive volcanic vents. In a joint model of hollow-hosting, degradation state and insolation account for the largest shares of the variation that physical factors explain. Hollow depth increases with hollow size, differs between host craters, varies little with temperature, host size, and composition, and does not change with nesting. Complete collapse of a charged layer whose base lies where sulfur vapor at 820--1,050~K balances the overburden produces the median depth of hollows, if sulfur and its pore space filled 10--40\% of the layer. The untested steps of the model (Section~\ref{sec:predictions}) can be tested with coupled models of heat and vapor transport in cooling craters, laboratory experiments, and BepiColombo images and thermal infrared spectra.

\section*{Acknowledgements}

I thank the MESSENGER team and the authors of the hollow, crater, and basin catalogs for releasing their data. I used Claude (Anthropic; Claude Opus models accessed through Claude Code and Claude Cowork, 2026) to write and run analysis code, to search for and check literature, to make figures, and to format and edit text. The code it wrote is archived with the other code of this study \parencite{zenodo_this}. I checked the code, results, and text and take responsibility for the content.

\section*{Open science declarations}

\noindent\textbf{Author contributions.} Michael S. Phillips: Conceptualization, Data curation, Formal analysis, Investigation, Methodology, Software, Validation, Visualization, Writing (original draft), and Writing (review \& editing).

\noindent\textbf{Code availability.} The analysis code (Python), including the georeferencing of hollow detections, the coverage and detection model, the thermal model, the crater containment graph, the terrain, slope, and depth analyses, the automatic shadow depths, the calculations of Section~S13, the color analysis of Section~S9, the joint model of Section~S16, and the figure scripts, is archived at Zenodo \parencite{zenodo_this} under the MIT License and is developed at https://github.com/Michael-S-Phillips/mercury-hollows. The crater containment graph and its interactive map are also archived separately \parencite{zenodo_cratergraph}. Image positions were computed with ISIS \parencite{isis2023}. Analyses used NumPy \parencite{harris2020}, SciPy \parencite{virtanen2020}, pandas \parencite{mckinney2010}, statsmodels \parencite{seabold2010}, and Matplotlib \parencite{hunter2007}.

\noindent\textbf{Data availability.} The hollow detections are from \textcite{bickel2025}, the crater catalog from \textcite{herrick2018}, the basin catalog from \textcite{orgel2020}, the tectonic map from \textcite{klimczak2025data}, the volcanic vents from \textcite{pegg2021}, the degradation stages from \textcite{deutsch2025}, the map of concentrated low-reflectance material from \textcite{klima2018}, and the hollow depths from \textcite{blewett2016}. MDIS images and the MDIS image index \parencite{mdis_cdr}, the MDIS eight-filter color map \parencite{mdis_mdr}, MESSENGER navigation data \parencite{spice_messenger}, and the X-Ray Spectrometer ratio maps \parencite{xrs_maps} are available from the NASA Planetary Data System. The global topography is available from the U.S. Geological Survey \parencite{usgs_dem}. The georeferenced hollow positions, my shadow-length depths (automatic and by hand), hollow fields, coverage and completeness grids, plains mask, terrain and slope classes, thermal-model grid, and the tables behind each figure are archived at Zenodo under a CC BY 4.0 license \parencite{zenodo_this}.

\noindent\textbf{Funding.} This work received no external funding.

\noindent\textbf{Competing interests.} The author serves on the editorial board of Planetary Research. The author declares no other known competing financial interests or personal relationships that could have appeared to influence the work reported in this paper.

% ------ Bibliography -------------------------------------
\printbibliography

\end{document}
